
TFHE supports homomorphic computation on encoded finite-domain integers, especially Boolean circuits. Its fast bootstrapping is useful for deep circuits when the encoding and noise bounds hold. Programmable bootstrapping combines ciphertext refresh with evaluation of a selected lookup function; this capability alone is not the formal cryptographic notion of functional encryption. The construction and error analysis are based on the original TFHE paper~\cite{tfhe-original}.

In TFHE, each plaintext is encrypted as an LWE ciphertext. Therefore, TFHE's ciphertext-to-ciphertext addition, ciphertext-to-plaintext addition, and ciphertext-to-plaintext multiplication are implemented based on GLWE's homomorphic addition and multiplication described in \autoref{part:generic-fhe}, with $n = 1$ to make GLWE an LWE.

This section will explain TFHE's novel components: key switching, ciphertext-to-ciphertext multiplication, coefficient extraction, and noise bootstrapping. 

$ $

\begin{tcolorbox}[
    title = \textbf{Required Background},    % box title
    colback = white,    % light background; tweak to taste
    colframe = black,  % frame colour
    boxrule = 0.8pt,     % line thickness
    left = 1mm, right = 1mm, top = 1mm, bottom = 1mm % inner padding
]
\begin{itemize}
  \item \autoref{sec:modulo}: \nameref{sec:modulo}
  \item \autoref{sec:group}: \nameref{sec:group}
  \item \autoref{sec:field}: \nameref{sec:field}
  \item \autoref{sec:order}: \nameref{sec:order}
  \item \autoref{sec:polynomial-ring}: \nameref{sec:polynomial-ring}
  \item \autoref{sec:decomp}: \nameref{sec:decomp}
  \item \autoref{sec:modulus-rescaling}: \nameref{sec:modulus-rescaling}
  \item \autoref{sec:lattice}: \nameref{sec:lattice}
  \item \autoref{sec:lwe}: \nameref{sec:lwe}
  \item \autoref{sec:rlwe}: \nameref{sec:rlwe}
  \item \autoref{sec:glwe}: \nameref{sec:glwe}
  \item \autoref{sec:glev}: \nameref{sec:glev}
  \item \autoref{sec:ggsw}: \nameref{sec:ggsw}
  \item \autoref{sec:glwe-add-cipher}: \nameref{sec:glwe-add-cipher}
  \item \autoref{sec:glwe-add-plain}: \nameref{sec:glwe-add-plain}
  \item \autoref{sec:glwe-mult-plain}: \nameref{sec:glwe-mult-plain}
  \item \autoref{subsec:modulus-switch-lwe}: \nameref{subsec:modulus-switch-lwe}
  \item \autoref{sec:glwe-key-switching}: \nameref{sec:glwe-key-switching}
\end{itemize}
\end{tcolorbox}

\clearpage

\subsection{Encryption and Decryption}
\label{subsec:tfhe-enc-dec}

TFHE encrypts and decrypts ciphertexts using LWE (\autoref{sec:lwe}), equivalently GLWE (\autoref{sec:glwe}) with $n=1$. The TFHE construction considered here samples the secret from $\{0,1\}^k$. Binary and ternary secrets are both choices within the broader LWE framework; a ternary secret is not part of the definition of LWE.


\begin{tcolorbox}[title={\textbf{\tboxlabel{\ref*{subsec:tfhe-enc-dec}} TFHE Encryption and Decryption}}]
\textbf{\underline{Initial Setup}:} $\Delta = \dfrac{q}{t}$, $\vec{s} \xleftarrow{\$} \mathbb{Z}_2^k$  \textcolor{red}{ $\rhd$ where $t$ divides $q$, and each element of $\vec{s}$ is a 0-degree polynomial}

$ $

\par\noindent\rule{\textwidth}{0.4pt}

\textbf{\underline{Encryption Input}:} $m \in \mathbb{Z}_t$, $\vec{a} \xleftarrow{\$} \mathbb{Z}_q^k$, $e \xleftarrow{\chi_\sigma} \mathbb{Z}_q$ \textcolor{red}{ $\rhd$ each element of $\vec{a}$ is a 0-degree polynomial}
\begin{enumerate}
\item Scale up $m \longrightarrow \Delta \cdot m \text{ } \in \mathbb{Z}_q$

\item Compute $b = \vec{a} \cdot \vec{s} + \Delta  m + e \text{ } \bmod q$
\item $\textsf{LWE}_{\vec{s},\sigma}(\Delta  m + e) = (\vec{a}, b) \text{ } \in \mathbb{Z}_q^{k + 1}$ 
\end{enumerate}

\par\noindent\rule{\textwidth}{0.4pt}

\textbf{\underline{Decryption Input}:} $\textsf{ct} = (\vec{a}, b) \text{ } \in \mathbb{Z}_q^{k+1}$
\begin{enumerate}
\item $\textsf{LWE}^{-1}_{\vec{s},\sigma}(\textsf{ct}) = b - \vec{a}\cdot \vec{s} = \Delta  m + e  \pmod q$

\item Scale down $\Bigg\lceil\dfrac{ \Delta  m + e } {\Delta}\Bigg\rfloor = m \text{ } \in \mathbb{Z}_t$ \textcolor{red}{ $\rhd$ i.e., modulus switch from $q \rightarrow t$}
\end{enumerate}

$ $

\textbf{{Condition for Correct Decryption}:}
\begin{itemize}
\item The noise $e$ grown over homomorphic operations should be: $|e| < \dfrac{\Delta}{2}$.
\end{itemize}

\end{tcolorbox}

In this section, we will often write $\textsf{LWE}_{\vec{s},\sigma}(\Delta  m + e)$ as $\textsf{LWE}_{\vec{s},\sigma}(\Delta  m)$ for simplicity, because $\textsf{LWE}_{\vec{s},\sigma}(\Delta  m + e) \approx \textsf{LWE}_{\vec{s},\sigma}(\Delta  m)$ (i.e., they decrypt to the same message). Even in the case that we write $\textsf{LWE}_{\vec{s},\sigma}(\Delta  m)$ instead of $\textsf{LWE}_{\vec{s},\sigma}(\Delta  m + e)$, you should assume this as an encryption of $\Delta m + e$ (i.e., the noise is included inside the scaled message).

\subsection{Homomorphic Ciphertext-to-Ciphertext Addition}
\label{subsec:tfhe-add-cipher}

TFHE's ciphertext-to-ciphertext addition uses LWE's ciphertext-to-ciphertext addition scheme, which is equivalent to GLWE's ciphertext-to-ciphertext addition scheme (\autoref{sec:glwe-add-cipher}) with $n = 1$.  


\begin{tcolorbox}[title={\textbf{\tboxlabel{\ref*{subsec:tfhe-add-cipher}} TFHE Ciphertext-to-Ciphertext Addition}}]
$\textsf{LWE}_{\vec{s}, \sigma}(\Delta m^{\langle 1 \rangle} + e^{\langle 1 \rangle} ) + \textsf{LWE}_{\vec{s}, \sigma}(\Delta m^{\langle 2 \rangle} + e^{\langle 2 \rangle}) $

$ = ( \vec{a}^{\langle 1 \rangle}, \text{ } b^{\langle 1 \rangle}) + (\vec{a}^{\langle 2 \rangle}, \text{ } b^{\langle 2 \rangle}) $

$ = ( \vec{a}^{\langle 1 \rangle} + \vec{a}^{\langle 2 \rangle}, \text{ } b^{\langle 1 \rangle} + b^{\langle 2 \rangle} ) $

$= \textsf{LWE}_{\vec{s}, \sigma}(\Delta(m^{\langle 1 \rangle} + m^{\langle 2 \rangle}) + e^{\langle 1 + 2 \rangle} )$
\label{Here}
\end{tcolorbox}


\subsection{Homomorphic Ciphertext-to-Plaintext Addition}
\label{subsec:tfhe-add-plain}

TFHE's ciphertext-to-plaintext addition (where $\lambda$ is a constant to add) uses LWE's ciphertext-to-plaintext addition scheme, which is equivalent to GLWE's ciphertext-to-plaintext addition scheme (\autoref{sec:glwe-add-plain}) with $n = 1$.  

\begin{tcolorbox}[title={\textbf{\tboxlabel{\ref*{subsec:tfhe-add-plain}} TFHE Ciphertext-to-Plaintext Addition}}]
$\textsf{LWE}_{\vec{s}, \sigma}(\Delta m + e) + \Delta \lambda $

$=  (\vec{a}, \text{ } b) + \Delta \lambda$

$=  (\vec{a}, \text{ } b + \Delta\lambda)$

$= \textsf{LWE}_{\vec{s}, \sigma}(\Delta (m + \lambda) + e )$
\end{tcolorbox}


\subsection{Homomorphic Ciphertext-to-Plaintext Multiplication}
\label{subsec:tfhe-mult-plain}

TFHE's ciphertext-to-plaintext multiplication uses LWE's ciphertext-to-plaintext multiplication scheme, which is equivalent to GLWE's ciphertext-to-plaintext multiplication scheme (\autoref{sec:glwe-mult-plain}) with $n = 1$.  

\begin{tcolorbox}[title={\textbf{\tboxlabel{\ref*{subsec:tfhe-mult-plain}} TFHE Ciphertext-to-Plaintext Multiplication}}]
$\textsf{LWE}_{\vec{s}, \sigma}(\Delta m + e) \cdot \lambda$

$= (\vec{a}, \text{ } b) \cdot \lambda$

$= (\lambda\cdot \vec{a}, \text{ } \lambda \cdot b)$

$= \textsf{LWE}_{\vec{s}, \sigma}(\Delta (m \cdot \lambda) + \lambda \cdot e )$
\end{tcolorbox}


\subsection{Homomorphic Key Switching}
\label{subsec:tfhe-key-switching}

\textbf{- Reference:} 
\href{https://www.zama.ai/post/tfhe-deep-dive-part-3}{TFHE Deep Dive - Part III - Key switching and leveled multiplications}~\cite{tfhe-3}

TFHE key switching changes the LWE secret from $\vec{s}$ to $\vec{s}_{'}$, possibly changing its dimension. With this chapter's phase $b-\langle\vec a,\vec s\rangle$, it is the $n=1$ instance of \autoref{sec:glwe-key-switching}. The minus sign below corresponds to this convention; the alternative convention in \autoref{subsec:glwe-alternative} uses the opposite sign.

\begin{tcolorbox}[title={\textbf{\tboxlabel{\ref*{subsec:tfhe-key-switching}} TFHE Key Switching}}]

Given $\textsf{LWE}_{\vec{s},\sigma}(\Delta m + e) = (\vec{a}, b)$,

$\textsf{LWE}_{\vec{s}_{'},\sigma}(\Delta m + e') = (0, b) - \bm{\langle} \textsf{Decomp}^{\beta, l}(\vec{a}), \text{ } \textsf{Lev}_{\vec{s}_{'}, \sigma}^{\beta, l}(\vec{s}) \bm{\rangle}$
\end{tcolorbox}


\subsection{Homomorphic Ciphertext-to-Ciphertext Multiplication}
\label{subsec:tfhe-mult-cipher}

The exact decomposition identities in the following proof assume $q=\beta^l$. For a truncated decomposition write $c=\langle\textsf{Decomp}(c),\vec g\rangle+\rho_c\pmod q$. Its residual must be included in the output error, together with the errors in the GSW rows. The notation $\sigma$ in a homomorphic output does not assert that its noise distribution remains the fresh encryption distribution.

\textbf{- Reference:} 
\href{https://www.zama.ai/post/tfhe-deep-dive-part-3}{TFHE Deep Dive - Part III - Key switching and leveled multiplications}~\cite{tfhe-3}

$ $

TFHE supports multiplication of two ciphertexts in the form: $\textsf{LWE}_{\vec{s}, \sigma}(\Delta m_1) \cdot \textsf{GSW}_{\vec{s}, \sigma}^{\beta, l}(m_2)$. 

$ $

\noindent The 1st term $\textsf{LWE}_{\vec{s}, \sigma}(\Delta m_1 + e_1)$ comes from one of the following: 
\begin{itemize}
\item A fresh LWE encryption (\autoref{subsec:glwe-enc}) of plaintext $m_1$. 
\item A homomorphically added result of two LWE ciphertexts (\autoref{sec:glwe-add-cipher}). 
\item A homomorphically multiplied result of a LWE ciphertext with a plaintext (\autoref{sec:glwe-mult-plain}). 
\end{itemize}

$ $

\noindent The 2nd term $\textsf{GSW}_{\vec{s}, \sigma}^{\beta, l}(m_2)$ comes from one of the following:
\begin{itemize}
\item A fresh GSW encryption (\autoref{subsec:ggsw-enc}) of plaintext $m_2$.
\item Converted from $\textsf{LWE}_{\vec{s}, \sigma}(\Delta m_2 + e_2)$ into $\textsf{GSW}_{\vec{s}, \sigma}^{\beta, l}(m_2)$ by \textit{circuit bootstrapping} (this will be covered in the future).
\end{itemize}

$ $

\noindent Remember the following: 

$\textsf{LWE}_{\vec{s}, \sigma}(\Delta m_1 + e_1) = (\vec{a}, b) \in \mathbb{Z}_{q}^{k + 1}$, where $b = \vec{a} \cdot \vec{s} + \Delta m_1 + e_1$ 

$ $

$\textsf{GSW}_{\vec{s}, \sigma}^{\beta, l}(m_2) = \Bigr \{ \{ \textsf{Lev}_{\vec{s}, \sigma}^{\beta, l} (-s_i \cdot m_2)  \}_{i=0}^{k-1}, \textsf{Lev}_{\vec{s}, \sigma}^{\beta, l}(m_2) \Bigl \} \in \mathbb{Z}_{q }^{(k+1) \cdot l \cdot (k+1)}$ \textcolor{red}{ $\rhd$ from \autoref{subsec:ggsw-enc}}

$ $

\noindent Let's use the following notations:

$\textsf{GSW}_{\vec{s}, \sigma}^{\beta, l}(m_2) = {\bar{\textsf{ct}}} = (\bar{\textsf{ct}}_0,. \bar{\textsf{ct}}_1, \gap{$\cdots$} \bar{\textsf{ct}}_k)$ 

$\bar{\textsf{ct}}_i = \textsf{Lev}_{\vec{s}, \sigma}^{\beta, l}(-s_i \cdot m_2)$ for $0 \leq i \leq (k-1)$

$\bar{\textsf{ct}}_k = \textsf{Lev}_{\vec{s}, \sigma}^{\beta, l}(m_2)$

$\textsf{ct} = \textsf{LWE}_{\vec{s}, \sigma}(\Delta m_1 + e_1) = (\vec{a}, b) = (a_0, a_1, \gap{$\cdots$}, a_{k-1}, b) = (\textsf{ct}_0, \textsf{ct}_1, \cdots, \textsf{ct}_k)$

$ $

\noindent Let's define the following TFHE ciphertext multiplication operation: 

$\textsf{ct} \cdot {\bar{\textsf{ct}}} = \sum\limits_{i=0}^{k}\langle \textsf{Decomp}^{\beta, l}(\textsf{ct}_i), \bar{\textsf{ct}}_i \rangle$

$ $

\noindent Then, the following is true:

\begin{tcolorbox}[title={\textbf{\tboxlabel{\ref*{subsec:tfhe-mult-cipher}} TFHE Ciphertext-to-Ciphertext Multiplication}}]
$\textsf{ct} = \textsf{LWE}_{\vec{s}, \sigma}(\Delta m_1 + e_1) = (a_0, a_1, \cdots, a_{k-1}, b)$

$\bar{\textsf{ct}} = \textsf{GSW}_{\vec{s}, \sigma}^{\beta, l}(m_2) = \bm( \textsf{Lev}_{\vec{s}, \sigma}^{\beta, l}(-s_0\cdot m_2), \textsf{Lev}_{\vec{s}, \sigma}^{\beta, l}(-s_1\cdot m_2), \cdots, \textsf{Lev}_{\vec{s}, \sigma}^{\beta, l}(-s_{k-1}\cdot m_2), \textsf{Lev}_{\vec{s}, \sigma}^{\beta, l}(m_2)  \bm)$

$\textsf{LWE}_{\vec{s}, \sigma}(\Delta m_1 + e_1) \cdot \textsf{GSW}_{\vec{s}, \sigma}^{\beta, l}(m_2) $

$= \sum\limits_{i=0}^{k}\langle \textsf{Decomp}^{\beta, l}(\textsf{ct}_i), \bar{\textsf{ct}}_i \rangle $

$\approx \textsf{LWE}_{\vec{s}, \sigma}(\Delta m_1 m_2)$
\end{tcolorbox}

This means that multiplying two TFHE ciphertexts (one is in LWE and another in GSW) and decrypting the resulting LWE ciphertext gives the same result as multiplying their two original plaintexts. 


%\textbf{\underline{Proof}}
\begin{myproof}
\begin{enumerate}
\item $\sum\limits_{i=0}^{k}\langle \textsf{Decomp}^{\beta, l}(\textsf{ct}_i), \bar{\textsf{ct}}_i \rangle$ \\
$= \langle \textsf{Decomp}^{\beta, l}(a_0), \bar{\textsf{ct}}_0 \rangle + \langle \textsf{Decomp}^{\beta, l}(a_1), \bar{\textsf{ct}}_1 \rangle + \gap{$\cdots$} + \langle \textsf{Decomp}^{\beta, l}(a_{k-1}), \bar{\textsf{ct}}_{k-1} \rangle + \langle \textsf{Decomp}^{\beta, l}(b), \bar{\textsf{ct}}_k \rangle$ \\
\textcolor{red}{ $\rhd$ expanding the dot product of two vectors}
\item For $i = k$: \\
$\textsf{Decomp}^{\beta, l}(b) = (b_1, b_2, \ldots , b_l)$, where $b = b_1\dfrac{q}{\beta^1} + b_2\dfrac{q}{\beta^2} + \cdots + b_l\dfrac{q}{\beta^l}$ \textcolor{red}{ $\rhd$ from \autoref{subsec:poly-decomp}}\\
$\bar{\textsf{ct}}_k = \textsf{Lev}_{\vec{s}, \sigma}^{\beta, l}(m_2) = \left(\textsf{LWE}_{\vec{s}, \sigma}\left(m_{2}\dfrac{q}{\beta^1} + e_{2, 1}\right), \textsf{LWE}_{\vec{s}, \sigma}\left(m_{2}\dfrac{q}{\beta^2} + e_{2, 2}\right), \ldots, \textsf{LWE}_{\vec{s}, \sigma}\left(m_{2}\dfrac{q}{\beta^l} + e_{2, l}\right) \right)$ \\
$ $

Therefore: \\
$\langle \textsf{Decomp}^{\beta, l}(b), \bar{\textsf{ct}}_k \rangle$\\
$= b_1 \cdot \textsf{LWE}_{\vec{s}, \sigma} \left (m_{2}\dfrac{q}{\beta^1}  + e_{2, 1}\right ) + b_2 \cdot \textsf{LWE}_{\vec{s}, \sigma} \left (m_{2}\dfrac{q}{\beta^2}  + e_{2, 2}\right ) + \gap{$\cdots$} + b_l \cdot \textsf{LWE}_{\vec{s}, \sigma} \left (m_{2}\dfrac{q}{\beta^l}  + e_{2, l}\right)$ \\
$= \textsf{LWE}_{\vec{s}, \sigma} \left (b_1m_{2}\dfrac{q}{\beta^1}  + b_1e_{2, 1}\right ) + \textsf{LWE}_{\vec{s}, \sigma} \left (b_2m_{2}\dfrac{q}{\beta^2}  + b_2e_{2,2}\right ) + \gap{$\cdots$} + \textsf{LWE}_{\vec{s}, \sigma} \left (b_lm_{2}\dfrac{q}{\beta^l} + b_le_{2, l}\right)$ \textcolor{red}{ $\rhd$ from \autoref{sec:glwe-mult-plain}} \\
$= \textsf{LWE}_{\vec{s}, \sigma} \left (b_1m_{2}\dfrac{q}{\beta^1} + b_2m_{2}\dfrac{q}{\beta^2} + \gap{$\cdots$} + b_lm_{2}\dfrac{q}{\beta^l} + b_1e_{2,1} + \cdots + b_le_{2,l}\right)$ \textcolor{red}{ $\rhd$ from \autoref{sec:glwe-add-cipher}}\\
$= \textsf{LWE}_{\vec{s}, \sigma} \left (m_{2} \cdot \left ( b_1\dfrac{q}{\beta^1} + b_2\dfrac{q}{\beta^2} + \gap{$\cdots$} + b_l\dfrac{q}{\beta^l} \right)+e^{\langle k \rangle}\right)$ \textcolor{red}{ $\rhd$ where $e^{\langle k \rangle} = \sum\limits_{i=1}^{l}b_ie_{2, i}$}\\
$= \textsf{LWE}_{\vec{s}, \sigma} (m_{2}b + e^{\langle k \rangle})$ \textcolor{red}{ $\rhd$ from \autoref{subsec:poly-decomp}}

\item For $0 \leq i \leq (k - 1)$: \\
$\textsf{Decomp}^{\beta, l}(a_i) = (a_{\langle i, 1 \rangle}, a_{\langle i, 2 \rangle}, \gap{$\cdots$}, a_{\langle i, l \rangle})$, where $a_i = a_{\langle i, 1 \rangle}\dfrac{q}{\beta^1} + a_{\langle i, 2 \rangle}\dfrac{q}{\beta^2} + \gap{$\cdots$} + a_{\langle i, l \rangle}\dfrac{q}{\beta^l}$ \\
$\bar{\textsf{ct}}_i =  \textsf{Lev}_{\vec{s}, \sigma}^{\beta, l}(-s_im_2) $

$= \left(\textsf{LWE}_{\vec{s}, \sigma}\left(-s_im_2\dfrac{q}{\beta^1} + e_{i, 1}\right), \textsf{LWE}_{\vec{s}, \sigma}\left(-s_im_2\dfrac{q}{\beta^2} + e_{i, 2}\right), \gap{$\cdots$}, \textsf{LWE}_{\vec{s}, \sigma}\left(-s_im_2\dfrac{q}{\beta^l} + e_{i, l}\right) \right)$ \\

$ $
Therefore: \\
$\langle \textsf{Decomp}^{\beta, l}(a_0), \bar{\textsf{ct}}_0 \rangle + \langle \textsf{Decomp}^{\beta, l}(a_1), \bar{\textsf{ct}}_1 \rangle + \gap{$\cdots$} + \langle \textsf{Decomp}^{\beta, l}(a_{k-1}), \bar{\textsf{ct}}_{k-1} \rangle$ \\
$= \sum\limits_{i=0}^{k-1}\langle \textsf{Decomp}^{\beta, l}(a_i), \bar{\textsf{ct}}_i \rangle$ \\
$= \sum\limits_{i=0}^{k-1}\Bigg(a_{\langle i, 1\rangle} \cdot \textsf{LWE}_{\vec{s}, \sigma}\left(-s_im_2\dfrac{q}{\beta^1} + e_{i,1}\right) + a_{\langle i, 2\rangle} \cdot \textsf{LWE}_{\vec{s}, \sigma}\left(-s_im_2\dfrac{q}{\beta^2} + e_{i,2}\right) + $

$ \cdots +  a_{\langle i, l\rangle} \cdot \textsf{LWE}_{\vec{s}, \sigma}\left(-s_im_2\dfrac{q}{\beta^l} + e_{i,l}\right)\Bigg)$ \\
$= \sum\limits_{i=0}^{k-1}\Bigg(\textsf{LWE}_{\vec{s}, \sigma}\left(-a_{\langle i, 1\rangle}s_im_2\dfrac{q}{\beta^1} + a_{\langle i, 1 \rangle}e_{i,1}\right) + \textsf{LWE}_{\vec{s}, \sigma}\left(-a_{\langle i, 2\rangle}s_im_2\dfrac{q}{\beta^2} + a_{\langle i, 2 \rangle}e_{i,2}\right) + $

$\gap{$\cdots$} + \textsf{LWE}_{\vec{s}, \sigma}\left(-a_{\langle i, l\rangle}s_im_2\dfrac{q}{\beta^l} + a_{\langle i, l \rangle}e_{i,l}\right)\Bigg)$ \\
$= \sum\limits_{i=0}^{k-1}\textsf{LWE}_{\vec{s}, \sigma}\left(-a_{\langle i, 1\rangle}s_im_2\dfrac{q}{\beta^1} -a_{\langle i, 2\rangle}s_im_2\dfrac{q}{\beta^2} + \gap{$\cdots$} -a_{\langle i, l\rangle}s_im_2\dfrac{q}{\beta^l} + a_{\langle i, 1 \rangle}e_{i,1} + \cdots + a_{\langle i, l \rangle}e_{i,l}\right)$ \\
$= \sum\limits_{i=0}^{k-1}\textsf{LWE}_{\vec{s}, \sigma}\left(-s_im_2 \cdot \left(a_{\langle i, 1\rangle}\dfrac{q}{\beta^1} + a_{\langle i, 2\rangle}\dfrac{q}{\beta^2} + \gap{$\cdots$} + a_{\langle i, l\rangle}\dfrac{q}{\beta^l}\right) + a_{\langle i, 1 \rangle}e_{i,1} + \cdots + a_{\langle i, l \rangle}e_{i,l}\right)$ \\
$= \sum\limits_{i=0}^{k-1}\textsf{LWE}_{\vec{s}, \sigma}(-s_im_2a_i + e^{\langle i \rangle})$ \textcolor{red}{ $\rhd$ where $e^{\langle i \rangle} = a_{\langle i, 1 \rangle}e_{i,1} + \cdots + a_{\langle i, l \rangle}e_{i,l}$}
\item According to step 2 and 3, \\
$\sum\limits_{i=0}^{k}\langle \textsf{Decomp}^{\beta, l}(\textsf{ct}_i), \bar{\textsf{ct}}_i \rangle$ \\
$= \sum\limits_{i=0}^{k-1}\textsf{LWE}_{\vec{s}, \sigma}(-s_im_2a_i + e^{\langle i \rangle}) + \textsf{LWE}_{\vec{s}, \sigma} (m_{2}b + e^{\langle k \rangle})$ \\ 
$= \textsf{LWE}_{\vec{s}, \sigma}\left(\sum\limits_{i=0}^{k-1}(-s_im_2a_i) + m_{2}b + \sum\limits_{i=0}^ke^{\langle i \rangle}\right)$ \textcolor{red}{ $\rhd$ addition of two LWE ciphertexts} \\ 
$= \textsf{LWE}_{\vec{s}, \sigma}\left(m_2b - \sum\limits_{i=0}^{k-1}m_2a_is_i + \sum\limits_{i=0}^ke^{\langle i \rangle}\right)$  \\ 
$= \textsf{LWE}_{\vec{s}, \sigma}\left(m_2(b - \sum\limits_{i=0}^{k-1}a_is_i) + \sum\limits_{i=0}^ke^{\langle i \rangle}\right)$\\
$= \textsf{LWE}_{\vec{s}, \sigma}\left(m_2(\Delta m_1 + e_1) + \sum\limits_{i=0}^ke^{\langle i \rangle}\right)$ \\
$= \textsf{LWE}_{\vec{s}, \sigma}\left(\Delta m_1m_2 + m_2e_1 + \sum\limits_{i=0}^ke^{\langle i \rangle}\right)$ \\
$\approx \textsf{LWE}_{\vec{s}, \sigma}(\Delta m_1m_2)$ \textcolor{red}{ $\rhd$ provided the complete error $m_2e_1+\sum\limits_{i=0}^ke^{\langle i\rangle}$ lies within the output decoding interval}
\end{enumerate}

\end{myproof}

%This issue of noise growth is similar to that in ciphertext-to-plaintext multiplication (\autoref{subsubsec:glwe-mult-plain-discussion}). 
To reduce the noise growth, noise bootstrapping is needed (will be discussed in \autoref{subsec:tfhe-noise-bootstrapping}).

\subsubsection{Generalization to GLWE-to-GGSW Multiplication}
\label{subsubsec:tfhe-glwe-to-ggsw-multiplication}

As in the scalar proof, exact descending gadget weights require $q=\beta^l$. The displayed calculation suppresses the GGSW row errors only to explain the message algebra. In full, if $\rho_{A_i},\rho_B$ are decomposition residuals and $E_{i,j}^{\rm key}$ are evaluation-key errors, the output error is
\[
E_{\rm out}=M_2E-M_2\Bigl(\rho_B-\sum_{i=0}^{k-1}\rho_{A_i}S_i\Bigr)+\sum_{i=0}^{k}\sum_{j=1}^{l}D_{i,j}E_{i,j}^{\rm key}\pmod q.
\]
All these terms require bounds. In particular, $\|M_2E\|_\infty\leq\|M_2\|_1\|E\|_\infty$; small input noise alone does not imply small output noise.

We can further generalize TFHE's LWE-to-GSW multiplication to GLWE-to-GGSW multiplication between the following two ciphertexts: $\textsf{GLWE}_{\vec{S}, \sigma}(\Delta M_1) \cdot \textsf{GGSW}_{\vec{S}, \sigma}^{\beta, l}(M_2)$, where $M_1$, $M_2$, and $S$ are $(n-1)$-degree polynomials. 


$ $

\noindent The 1st term $\textsf{GLWE}_{\vec{S}, \sigma}(\Delta M_1)$ comes from one of the following: 
\begin{itemize}
\item A fresh GLWE encryption (\autoref{subsec:glwe-enc}) of plaintext $M_1$. 
\item A homomorphically added result of two GLWE ciphertexts (\autoref{sec:glwe-add-cipher}). 
\item A homomorphically multiplied result of a GLWE ciphertext with a plaintext (\autoref{sec:glwe-mult-plain}). 
\end{itemize}

$ $

\noindent The 2nd term $\textsf{GGSW}_{\vec{S}, \sigma}^{\beta, l}(M_2)$ comes from one of the following:
\begin{itemize}
\item A fresh GGSW encryption (\autoref{subsec:ggsw-enc}) of plaintext $M_2$.
\item Converted from $\textsf{GLWE}_{\vec{S}, \sigma}(\Delta M_2)$ into $\textsf{GGSW}_{\vec{S}, \sigma}^{\beta, l}(M_2)$ by \textit{circuit bootstrapping} (this will be covered in the future).
\end{itemize}

$ $

\noindent Remember the following: 

$\textsf{GLWE}_{\vec{S}, \sigma}(\Delta M_1) = (A_0, A_1, \gap{$\cdots$}, A_{k-1}, B) \in \mathcal{R}_{n, q}^{k + 1}$, where $B = \sum\limits_{i=0}^{k-1}(A_i \cdot S_i) + \Delta M_1 + E$ \\ \textcolor{red}{ $\rhd$ from \autoref{subsec:glwe-enc}} 

$ $

$\textsf{GGSW}_{\vec{S}, \sigma}^{\beta, l}(M_2) = \Bigr \{ \{ \textsf{GLev}_{\vec{S}, \sigma}^{\beta, l} (-S_i \cdot M_2)  \}_{i=0}^{k-1}, \textsf{GLev}_{\vec{S}, \sigma}^{\beta, l}(M_2) \Bigl \} \in \mathcal{R}_{\langle n, q \rangle }^{(k+1) \cdot l \cdot (k+1)}$ \textcolor{red}{ $\rhd$ from \autoref{subsec:ggsw-enc}}

$ $

\noindent Let's use the following notations:

$\textsf{GGSW}_{\vec{S}, \sigma}^{\beta, l}(M_2) = {\bar{C}} = (\bar{C_0},. \bar{C_1}, \gap{$\cdots$} \bar{C_k})$ 

$\bar{C_i} = \textsf{GLev}_{\vec{S}, \sigma}^{\beta, l}(-S_i \cdot M_2)$ for $0 \leq i \leq (k-1)$

$\bar{C_k} = \textsf{GLev}_{\vec{S}, \sigma}^{\beta, l}(M_2)$

$\textsf{ct} = \textsf{GLWE}_{\vec{S}, \sigma}(\Delta M_1) = (C_0, C_1, \gap{$\cdots$}, C_k) = (A_0, A_1, \gap{$\cdots$}, A_{k-1}, B)$

$ $

\noindent Let's define the following TFHE ciphertext multiplication operation: 

$\textsf{ct} \cdot {\bar{C}} = \sum\limits_{i=0}^{k}\langle \textsf{Decomp}^{\beta, l}(C_i), \bar{C_i} \rangle$

$ $

\noindent Then, the following is true:

\begin{tcolorbox}[title={\textbf{\tboxlabel{\ref*{subsubsec:tfhe-glwe-to-ggsw-multiplication}} Generalization to GLWE-to-GGSW Multiplication}}]
$\textsf{ct} = \textsf{GLWE}_{\vec{S}, \sigma}(\Delta M_1) = (A_0, A_1, \cdots, A_{k-1}, B)$

$\bar{C} = \textsf{GGSW}_{\vec{S}, \sigma}^{\beta, l}(M_2)$

$ = \bm(\textsf{GLev}_{\vec{S}, \sigma}^{\beta, l}(-S_0\cdot M_2), \textsf{GLev}_{\vec{S}, \sigma}^{\beta, l}(-S_1\cdot M_2), \cdots, \textsf{GLev}_{\vec{S}, \sigma}^{\beta, l}(-S_{k-1}\cdot M_2), \textsf{GLev}_{\vec{S}, \sigma}^{\beta, l}(M_2)\bm)$

$ $

$\textsf{GLWE}_{\vec{S}, \sigma}(\Delta M_1) \cdot \textsf{GGSW}_{\vec{S}, \sigma}^{\beta, l}(M_2) = \sum\limits_{i=0}^{k}\langle \textsf{Decomp}^{\beta, l}(C_i), \bar{C_i} \rangle \approx \textsf{GLWE}_{\vec{S}, \sigma}(\Delta M_1 M_2)$
\end{tcolorbox}

This means that multiplying two TFHE ciphertexts (one is in GLWE and another in GGSW) and decrypting the resulting GLWE ciphertext gives the same result as multiplying their two original plaintexts. 


%\textbf{\underline{Proof}}
\begin{myproof}
\begin{enumerate}
\item $\sum\limits_{i=0}^{k}\langle \textsf{Decomp}^{\beta, l}(C_i), \bar{C_i} \rangle$ \\
$= \langle \textsf{Decomp}^{\beta, l}(A_0), \bar{C_0} \rangle + \langle \textsf{Decomp}^{\beta, l}(A_1), \bar{C_1} \rangle + \gap{$\cdots$} + \langle \textsf{Decomp}^{\beta, l}(A_{k-1}), \bar{C}_{k-1} \rangle + \langle \textsf{Decomp}^{\beta, l}(B), \bar{C_k} \rangle$ \\
\textcolor{red}{ $\rhd$ expanding the dot product of two vectors}
\item For $i = k$: \\
$\textsf{Decomp}^{\beta, l}(B) = (B_1, B_2, \gap{$\cdots$}, B_l)$, where $B = B_1\dfrac{q}{\beta^1} + B_2\dfrac{q}{\beta^2} + \gap{$\cdots$} + B_l\dfrac{q}{\beta^l}$ \textcolor{red}{ $\rhd$ from \autoref{subsec:poly-decomp}}\\
$\bar{C}_k = \textsf{GLev}_{\vec{S}, \sigma}^{\beta, l}(M_2) = \left(\textsf{GLWE}_{\vec{S}, \sigma}\left(M_{2}\dfrac{q}{\beta^1}\right), \textsf{GLWE}_{\vec{S}, \sigma}\left(M_{2}\dfrac{q}{\beta^2}\right), \gap{$\cdots$}, \textsf{GLWE}_{\vec{S}, \sigma}\left(M_{2}\dfrac{q}{\beta^l}\right) \right)$ 

\textcolor{red}{ $\rhd$ we omit the noise terms $E_{2,1}, \ldots, E_{2,l}$ in each GLWE ciphertext for simplicity} \\
$ $

Therefore: \\
$\langle \textsf{Decomp}^{\beta, l}(B), \bar{C_k} \rangle$\\
$= B_1 \cdot \textsf{GLWE}_{\vec{S}, \sigma} \left (M_{2}\dfrac{q}{\beta^1} \right ) + B_2 \cdot \textsf{GLWE}_{\vec{S}, \sigma} \left (M_{2}\dfrac{q}{\beta^2} \right ) + \gap{$\cdots$} + B_l \cdot \textsf{GLWE}_{\vec{S}, \sigma} \left (M_{2}\dfrac{q}{\beta^l}\right)$ \\
$= \textsf{GLWE}_{\vec{S}, \sigma} \left (B_1M_{2}\dfrac{q}{\beta^1} \right ) + \textsf{GLWE}_{\vec{S}, \sigma} \left (B_2M_{2}\dfrac{q}{\beta^2} \right ) + \gap{$\cdots$} + \textsf{GLWE}_{\vec{S}, \sigma} \left (B_lM_{2}\dfrac{q}{\beta^l}\right)$ \textcolor{red}{ $\rhd$ from \autoref{sec:glwe-mult-plain}} \\
$= \textsf{GLWE}_{\vec{S}, \sigma} \left (B_1M_{2}\dfrac{q}{\beta^1} + B_2M_{2}\dfrac{q}{\beta^2} + \gap{$\cdots$} + B_lM_{2}\dfrac{q}{\beta^l}\right)$ \textcolor{red}{ $\rhd$ from \autoref{sec:glwe-add-cipher}}\\
$= \textsf{GLWE}_{\vec{S}, \sigma} \left (M_{2} \cdot \left ( B_1\dfrac{q}{\beta^1} + B_2\dfrac{q}{\beta^2} + \gap{$\cdots$} + B_l\dfrac{q}{\beta^l} \right)\right)$ \\
$= \textsf{GLWE}_{\vec{S}, \sigma} (M_{2}B)$ \textcolor{red}{ $\rhd$ from \autoref{subsec:poly-decomp}}
\item For $0 \leq i \leq (k - 1)$: \\
$\textsf{Decomp}^{\beta, l}(A_i) = (A_{\langle i, 1 \rangle}, A_{\langle i, 2 \rangle}, \gap{$\cdots$}, A_{\langle i, l \rangle})$, where $A_i = A_{\langle i, 1 \rangle}\dfrac{q}{\beta^1} + A_{\langle i, 2 \rangle}\dfrac{q}{\beta^2} + \gap{$\cdots$} + A_{\langle i, l \rangle}\dfrac{q}{\beta^l}$ \\
$\bar{C_i} =  \textsf{GLev}_{\vec{S}, \sigma}^{\beta, l}(-S_iM_2) = \left(\textsf{GLWE}_{\vec{S}, \sigma}\left(-S_iM_2\dfrac{q}{\beta^1}\right), \textsf{GLWE}_{\vec{S}, \sigma}\left(-S_iM_2\dfrac{q}{\beta^2}\right), \gap{$\cdots$}, \textsf{GLWE}_{\vec{S}, \sigma}\left(-S_iM_2\dfrac{q}{\beta^l}\right) \right)$ \\

$ $
Therefore: \\
$\langle \textsf{Decomp}^{\beta, l}(A_0), \bar{C_0} \rangle + \langle \textsf{Decomp}^{\beta, l}(A_1), \bar{C_1} \rangle + \gap{$\cdots$} + \langle \textsf{Decomp}^{\beta, l}(A_{k-1}), \bar{C}_{k-1} \rangle$ \\
$= \sum\limits_{i=0}^{k-1}\langle \textsf{Decomp}^{\beta, l}(A_i), \bar{C_i} \rangle$ \\
$= \sum\limits_{i=0}^{k-1}\left(A_{\langle i, 1\rangle} \cdot \textsf{GLWE}_{\vec{S}, \sigma}\left(-S_iM_2\dfrac{q}{\beta^1}\right) + A_{\langle i, 2\rangle} \cdot \textsf{GLWE}_{\vec{S}, \sigma}\left(-S_iM_2\dfrac{q}{\beta^2}\right) + \gap{$\cdots$} + A_{\langle i, l\rangle} \cdot \textsf{GLWE}_{\vec{S}, \sigma}\left(-S_iM_2\dfrac{q}{\beta^l}\right)\right)$ \\
$= \sum\limits_{i=0}^{k-1}\left(\textsf{GLWE}_{\vec{S}, \sigma}\left(-A_{\langle i, 1\rangle}S_iM_2\dfrac{q}{\beta^1}\right) + \textsf{GLWE}_{\vec{S}, \sigma}\left(-A_{\langle i, 2\rangle}S_iM_2\dfrac{q}{\beta^2}\right) + \gap{$\cdots$} + \textsf{GLWE}_{\vec{S}, \sigma}\left(-A_{\langle i, l\rangle}S_iM_2\dfrac{q}{\beta^l}\right)\right)$ \\
$= \sum\limits_{i=0}^{k-1}\textsf{GLWE}_{\vec{S}, \sigma}\left(-A_{\langle i, 1\rangle}S_iM_2\dfrac{q}{\beta^1} + -A_{\langle i, 2\rangle}S_iM_2\dfrac{q}{\beta^2} + \gap{$\cdots$} + -A_{\langle i, l\rangle}S_iM_2\dfrac{q}{\beta^l}\right)$ \\
$= \sum\limits_{i=0}^{k-1}\textsf{GLWE}_{\vec{S}, \sigma}\left(-S_iM_2 \cdot \left(A_{\langle i, 1\rangle}\dfrac{q}{\beta^1} + A_{\langle i, 2\rangle}\dfrac{q}{\beta^2} + \gap{$\cdots$} + A_{\langle i, l\rangle}\dfrac{q}{\beta^l}\right)\right)$ \\
$= \sum\limits_{i=0}^{k-1}\textsf{GLWE}_{\vec{S}, \sigma}(-S_iM_2A_i)$
\item According to step 2 and 3, \\
$\sum\limits_{i=0}^{k}\langle \textsf{Decomp}^{\beta, l}(C_i), \bar{C_i} \rangle$ \\
$= \sum\limits_{i=0}^{k-1}\textsf{GLWE}_{\vec{S}, \sigma}(-S_iM_2A_i) + \textsf{GLWE}_{\vec{S}, \sigma} (M_{2}B)$ \\ 
$= \textsf{GLWE}_{\vec{S}, \sigma}\Big(\sum\limits_{i=0}^{k-1}(-S_iM_2A_i) + M_{2}B\Big)$ \textcolor{red}{ $\rhd$ addition of two GLWE ciphertexts} \\ 
$= \textsf{GLWE}_{\vec{S}, \sigma}\Big(BM_2 - \sum\limits_{i=0}^{k-1}M_2A_iS_i\Big)$  \\ 
$= \textsf{GLWE}_{\vec{S}, \sigma}\Big(M_2(B - \sum\limits_{i=0}^{k-1}A_iS_i)\Big)$\\
$= \textsf{GLWE}_{\vec{S}, \sigma}(M_2(\Delta M_1 + E))$ \\
$= \textsf{GLWE}_{\vec{S}, \sigma}(\Delta M_1M_2 + M_2E)$ \\
$\approx \textsf{GLWE}_{\vec{S}, \sigma}(\Delta M_1M_2)$ \textcolor{red}{ $\rhd$ provided the full output error satisfies the decoding bound}
\end{enumerate}

\end{myproof}

\subsection{Coefficient Extraction}
\label{subsec:tfhe-extraction}

\textbf{- Reference:} 
\href{https://www.zama.ai/post/tfhe-deep-dive-part-4}{TFHE Deep Dive - Part IV - Programmable Bootstrapping}~\cite{tfhe-4}

$ $

In TFHE, coefficient extraction is a process of extracting a coefficient of a polynomial that is encrypted as GLWE ciphertext. The extracted coefficient is in the form of LWE ciphertext (\autoref{sec:lwe}). %We will explain how to extract the coefficient of a plaintext polynomial $M$ from a GLWE ciphertext and RLWE ciphertext, respectively. 


%\subsubsection{Coefficient Extraction from a GLWE Ciphertext}
%\label{subsubsec:tfhe-extraction-glwe}

Note that in the GLWE cryptosystem, plaintext $M$ is encoded as a polynomial, where each coefficient encodes the plaintext value $m_0, m_1, \cdots, m_{n-1}$.

%\subsubsection{Overview}
%\label{subsec:tfhe-extraction-overview}

Suppose we have a GLWE ciphertext setup as the following: \\ 
$M = \sum\limits_{j=0}^{n-1}m_jX^j \in \mathcal{R}_{\langle n, t \rangle}$

$S = \left(S_0 = \sum\limits_{j=0}^{n-1}s_{0,j}X^j, S_1 = \sum\limits_{j=0}^{n-1}s_{1,j}X^j, \gap{$\cdots$}, S_{k-1} = \sum\limits_{j=0}^{n-1}s_{k-1,j}X^j \right)$ 

$\textsf{GLWE}_{\vec{S}, \sigma}(\Delta M) = \left(A_0 = \sum\limits_{j=0}^{n-1}a_{0,j}X^j, A_1 = \sum\limits_{j=0}^{n-1}a_{1,j}X^j, \gap{$\cdots$}, A_{k-1} = \sum\limits_{j=0}^{n-1}a_{k-1,j}X^j, B = \sum\limits_{j=0}^{n-1}b_{j}X^j\right)$ 

$B = \sum\limits_{i=0}^{k-1}A_iS_i + \Delta M + E$ 

$E = \sum\limits_{i=0}^{n-1}e_iX^i$ 

$ $

\noindent Note that:

$\Delta M + E = B - \sum\limits_{i=0}^{k-1}A_iS_i$ 

$ = (\Delta m_0 + \Delta m_1X + \gap{$\cdots$} + \Delta m_{n-1}X^{n-1}) + (e_0 + e_1X + \gap{$\cdots$} + e_{n-1}X^{n-1})$

$= (\Delta m_0 + e_0) + (\Delta m_1 + e_1)X + \gap{$\cdots$} + (\Delta m_{n-1} + e_{n-1})X^{n-1}$

$ $

\noindent Another way to write the formula is:

$B - \sum\limits_{i=0}^{k-1}A_iS_i$ 

$ = (b_0 + b_1X + \gap{$\cdots$} + b_{n-1}X^{n-1} )$ 

$ - (a_{0,0} + a_{0,1}X + \gap{$\cdots$} + a_{0, n-1}X^{n-1})(s_{0,0} + s_{0,1}X + \gap{$\cdots$} + s_{0, n-1}X^{n-1})$ 

$ - (a_{1,0} + a_{1,1}X + \gap{$\cdots$} + a_{1, n-1}X^{n-1})(s_{1,0} + s_{1,1}X + \gap{$\cdots$} + s_{1, n-1}X^{n-1})$

$ - \gap{$\cdots$} $ 

$ - (a_{k-1,0} + a_{k-1,1}X + \gap{$\cdots$} + a_{k-1, n-1}X^{n-1})(s_{k-1,0} + s_{k-1,1}X + \gap{$\cdots$} + s_{k-1, n-1}X^{n-1})$ 

$ $

$ = \left(b_0 - \left( \sum\limits_{i=0}^{k-1} \sum\limits_{j=0}^{0}(a_{i,0-j}s_{i,j}) - \sum\limits_{i=0}^{k-1} \sum\limits_{j=1}^{n-1}(a_{i,n-j}s_{i,j}) \right)\right)$

$ + \left(b_1 - \left( \sum\limits_{i=0}^{k-1} \sum\limits_{j=0}^{1}(a_{i,1-j}s_{i,j}) - \sum\limits_{i=0}^{k-1} \sum\limits_{j=2}^{n-1}(a_{i,n+1-j}s_{i,j})   \right) \right)\cdot X$ 



$ + \left(b_2 - \left( \sum\limits_{i=0}^{k-1} \sum\limits_{j=0}^{2}(a_{i,2-j}s_{i,j}) - \sum\limits_{i=0}^{k-1} \sum\limits_{j=3}^{n-1}(a_{i,n+2-j}s_{i,j})   \right) \right)\cdot X^2$ 

$ $

$\gap{$\cdots$}$ 

$ $

$ + \left(b_{n-1} - \left(  \sum\limits_{i=0}^{k-1} \sum\limits_{j=0}^{n-1}(a_{i,n-1-j}s_{i,j}) - \sum\limits_{i=0}^{k-1} \sum\limits_{j=n}^{n-1}(a_{i,n+(n-1)-j}s_{i,j})  \right) \right)\cdot X^{n-1}$ 

\textcolor{red}{ $\rhd$ Grouping the terms by same exponents}

$ $


$ $

$= \sum\limits_{h=0}^{n-1}  \left(b_h - \left(  \sum\limits_{i=0}^{k-1} \sum\limits_{j=0}^{h}(a_{i,h-j}s_{i,j}) - \sum\limits_{i=0}^{k-1} \sum\limits_{j=h+1}^{n-1}(a_{i,n+h-j}s_{i,j})  \right) \right)\cdot X^{h}  $


$ $

$= \sum\limits_{h=0}^{n-1}  C_h \cdot X^{h}  $, where $C_h = b_h - \left(  \sum\limits_{i=0}^{k-1} \sum\limits_{j=0}^{h}(a_{i,h-j}s_{i,j}) - \sum\limits_{i=0}^{k-1} \sum\limits_{j=h+1}^{n-1}(a_{i,n+h-j}s_{i,j})  \right)$

$ $

\noindent In the above $(n-1)$-degree polynomial, notice that each $X^h$ term's coefficient, $C_h$, can be expressed as an LWE ciphertext $\textsf{ct}_h$ as follows:

$S' = (s_{0,0}, s_{0,1}, \gap{$\cdots$}, s_{0,n-1}, s_{1,0}, s_{1,1}, \gap{$\cdots$}, s_{1, n-1}, \gap{$\cdots$}, s_{k-1, n-1}) = (s'_0, s'_1, \gap{$\cdots$}, s'_{nk-1} ) \in \mathbb{Z}_q^{nk}$


$\textsf{ct}_h = (a'_0, a'_1, \gap{$\cdots$}, a'_{nk-1}, b_h) \in \mathbb{Z}_q^{nk + 1}$


\[
    \text{, where } a'_{n \cdot i + j} =   
\begin{cases}
    a_{i,h - j} \text{ (if } 0 \leq j \leq h\text{)}\\
    -a_{i,n + h - j} \text{ (if } h+1 \leq j \leq n-1\text{)}\\
\end{cases}
\centering , b_h \text{ is directly obtained from the polynomial } B
\]


\noindent Note that $b_h - \sum\limits_{i=0}^{nk-1}s'_ia'_i = \Delta m_h + e_h$. This means that $C_h$ can be replaced by its encrypted version, $\textsf{LWE}_{\vec{s}_{'}, \sigma}(\Delta m_h)$, an LWE ciphertext $\textsf{ct}_h$ encrypting the $h$-th coefficient of $M$. Therefore, we just extracted $\textsf{LWE}_{\vec{s}_{'}, \sigma}(\Delta m_h)$ from $\textsf{GLWE}_{\vec{S}, \sigma}(\Delta M)$. This operation is called coefficient extraction, which does not add any noise because it simply extracts an LWE ciphertext by reordering the polynomial of the GLWE ciphertext. 

Once we have $\textsf{LWE}_{\vec{s}_{'}, \sigma}(\Delta m_h)$, we can key-switch it from $\vec{s}_{'} \rightarrow \vec{s}$ (\autoref{subsec:tfhe-key-switching}). 

\begin{tcolorbox}[title={\textbf{\tboxlabel{\ref*{subsec:tfhe-extraction}} GLWE Ciphertext's Coefficient Extraction}}]
Given the following GLWE ciphertext: 

$M = \sum\limits_{j=0}^{n-1}m_jX^j \in \mathcal{R}_{\langle n, t \rangle}$

$\vec{S} = \left(S_0 = \sum\limits_{j=0}^{n-1}s_{0,j}X^j, S_1 = \sum\limits_{j=0}^{n-1}s_{1,j}X^j, \gap{$\cdots$}, S_{k-1} = \sum\limits_{j=0}^{n-1}s_{k-1,j}X^j \right)$

$\textsf{GLWE}_{\vec{S}, \sigma}(\Delta M) = \left(A_0 = \sum\limits_{j=0}^{n-1}a_{0,j}X^j, A_1 = \sum\limits_{j=0}^{n-1}a_{1,j}X^j, \gap{$\cdots$}, A_{k-1} = \sum\limits_{j=0}^{n-1}a_{k-1,j}X^j, B = \sum\limits_{j=0}^{n-1}b_{j}X^j\right)$

$B = \sum\limits_{i=0}^{k-1}A_iS_i + \Delta M + E \bmod q$, \text{ } $E = \sum\limits_{i=0}^{n-1}e_iX^i$

$ $

$\textsf{LWE}_{\vec{s}_{'}, \sigma}(\Delta m_h)$ is an LWE ciphertext that encrypts $\Delta M$'s $h$-th coefficient (i.e., $\Delta m_h$). $\textsf{LWE}_{\vec{s}_{'}, \sigma}(\Delta m_h)$ can be extracted from $\textsf{GLWE}_{\vec{S}, \sigma}(\Delta M)$ as follows: 

$ $

$\vec{s}_{'} = (s_{0,0}, s_{0,1}, \gap{$\cdots$}, s_{0,n-1}, s_{1,0}, s_{1,1}, \gap{$\cdots$}, s_{1, n-1}, \gap{$\cdots$}, s_{k-1, n-1}) = (s'_0, s'_1, \gap{$\cdots$}, s'_{nk-1} ) \in \mathbb{Z}_q^{nk}$

$\textsf{LWE}_{\vec{s}_{'}, \sigma}(\Delta m_h) = (a_0', a_1', \gap{$\cdots$} , a_{nk-1}', b_h) \in \mathbb{Z}_q^{nk + 1}$


\[
    \text{, where } a'_{n \cdot i + j} =   
\begin{cases}
    a_{i,h - j} \text{ (if } 0 \leq j \leq h\text{)}\\
    -a_{i,n + h - j} \text{ (if } h+1 \leq j \leq n-1\text{)}\\
\end{cases}
, b_h \text{ is obtained from the polynomial } B
\]

Once we have $\textsf{LWE}_{\vec{s}_{'}, \sigma}(\Delta m_h)$, key-switch it from $\vec{s}_{'} \rightarrow \vec{s}$ (\autoref{subsec:tfhe-key-switching}).

\end{tcolorbox}



\subsection{Noise Bootstrapping}
\label{subsec:tfhe-noise-bootstrapping}


\textbf{- Reference:} 
\href{https://www.zama.ai/post/tfhe-deep-dive-part-4}{TFHE Deep Dive - Part IV - Programmable Bootstrapping}~\cite{tfhe-4}

$ $

Homomorphic addition generally increases the noise budget even when each fresh error has mean zero. For independent errors of variance $\sigma^2$, a sum of $r$ errors has variance $r\sigma^2$; adding the same ciphertext $r$ times gives error $re$ and variance $r^2\sigma^2$. Multiplication and external products also introduce message-dependent, evaluation-key and decomposition errors. Bootstrapping refreshes a ciphertext before its accumulated error crosses a decoding boundary (illustrated in \autoref{fig:scaling}); it replaces the old error by a new bounded evaluation error, rather than by zero.

As explained in the beginning of this section, TFHE uses LWE (which is GLWE with polynomial degree 0) to encrypt \& decrypt a plaintext. That is, each plaintext is $m$ (a single number), encoded as a zero-degree polynomial. Further, the secret key S that encrypts each $m$ is a vector $ \{s_0, s_1, \text{ } \cdots \text{ }, s_{k-1} \}$ instead of a polynomial. On the other hand, TFHE's noise bootstrapping uses homomorphic addition between GLWE ciphertexts and homomorphic multiplication between GLWE and GGSW ciphertexts.  

%Nonetheless, the reason why we learned TFHE ciphertext addition (\autoref{subsec:glwe-add} and \autoref{subsec:glwe-add-plain}), multiplication (\autoref{subsec:tfhe-mult-cipher} and \autoref{subsec:glwe-mult-plain}), and polynomial rotation (\autoref{subsec:coeff-rotation}) in terms of GLWE is because the bootstrapping procedure we will explain in this subsection requires homomorphic addition and multiplication of GLWE ciphertexts.

Suppose we have a TFHE ciphertext as follows: 

$ $

$\textsf{LWE}_{\vec{s}, \sigma}(\Delta m) = (a_0, a_1, \gap{$\cdots$} a_{k-1}, b)$

$b = \sum\limits_{i=0}^{k-1} a_is_i + \Delta m + e_b$

$\vec{s} = (s_0, s_1, \gap{$\cdots$} s_{k-1})$

$ $

, where $e_b$ is a big noise accumulated over a series of many ciphertext (or plaintext) multiplications. The goal of noise bootstrapping is to convert $(a_0, a_1, \gap{$\cdots$} a_{k-1}, b)$ into $(a_0', a_1', \gap{$\cdots$} a_{k-1}', b')$ such that: 

$ $

$b' = \sum\limits_{i=0}^{k-1} a_i's_i + \Delta m + e_s$

$ $

, where $e_s$ is a re-initialized noise. 

$ $
\subsubsection{Overview}
\label{subsec:bootstrapping-overview}

Bootstrapping evaluates a rounding lookup table at the encrypted phase. Conceptually, let $j=[\Delta m+e]_q$ and define a length-$q$ cyclic table by
\[
F_q(j)=\Bigl[\Bigl\lceil j/\Delta\Bigr\rfloor\Bigr]_t,
\qquad V_q(X)=\sum_{j=0}^{q-1}F_q(j)X^j.
\]
Here the centered error satisfies $|e|<\Delta/2$ and ties require a fixed convention. Thus each message occupies an interval centered at $\Delta m$, rather than the one-sided interval $[\Delta m,\Delta(m+1))$. This conceptual cyclic table explains lookup; it is not a GLWE ciphertext in the smaller ring $\mathcal R_{\langle n,q\rangle}$.

The actual construction uses a degree-at-most-$n-1$ polynomial $V(X)$, modulus switching to $2n$, and the negacyclic lookup restriction explained next. We use $V$ for the unscaled table and initialize a GLWE accumulator encrypting $\Delta V$. For a public table the trivial ciphertext $(0,\ldots,0,\Delta V)$ suffices. The output GLWE secret $\vec S_{bk}$ consists of $k'$ binary polynomials; it is secret and is not itself the public bootstrapping key. That key is the collection
\[
\textsf{BK}=\{\textsf{GGSW}_{\vec S_{bk},\sigma}^{\beta,l}(s_i)\}_{i=0}^{k-1}.
\]
Blind rotation uses these encryptions, not the secret $\vec S_{bk}$, to rotate the encrypted accumulator. Coefficient extraction (\autoref{subsec:tfhe-extraction}) then gives an LWE encryption of $\Delta m$ under the flattened GLWE key, followed by key switching back to $\vec s$.

\begin{enumerate}
\item \textbf{Input:} an LWE encryption of $\Delta m$ with error $e$.
\item Modulus-switch its public components to $2n$ using \autoref{subsec:modulus-switch-lwe}.
\item Blindly rotate the encryption of $\Delta V$ by the resulting encrypted phase, using the coefficient-rotation operation of Summary~\ref*{subsec:coeff-rotation}.1 in \autoref{subsec:coeff-rotation}.
\item Extract the constant coefficient and switch its secret key back to $\vec s$.
\item \textbf{Output:} an encryption of the same $m$, provided the input phase remains inside its lookup interval and the output evaluation error is below the decoding threshold.
\end{enumerate}

The input error affects which table entry is selected. A correct lookup makes the selected message independent of that error, while the homomorphic evaluation produces a new error. This distinction is essential: bootstrapping cannot recover a message after its input phase has crossed into another lookup interval.

\subsubsection{Modulus Switch for Noise Bootstrapping}
\label{subsec:bootstrapping-modulus-switch}

To avoid using the giant $(q-1)$-degree polynomial $V_q$, we will compress $q$ possible ciphertext elements $\Delta m + e \in \mathbb{Z}_q$ into $n$ distinct exponents of the $(n-1)$-degree polynomial $V$, where each $v_jX^j$ term in $V$ represents a mapping from $j \rightarrow v_j$ (i.e., noisy plaintext to noise-free plaintext). However, notice that when we rotate the coefficients of the $(n-1)$-degree polynomial $V$ to the left, as $v_jX^j$ rotates across the boundary between $X^0$ and $X^{n-1}$ degree terms, $v_j$'s sign flips to $-v_j$ (as shown in the example of \autoref{subsec:coeff-rotation-ex}). Due to this coefficient sign flip, the $(n-1)$-degree polynomial $V$ can theoretically encode total $2n$ distinct coefficient states as follows: $(v_0, v_1, v_2, \gap{$\cdots$}, v_{n-1}, -v_0, -v_1, \gap{$\cdots$}, -v_{n-1})$. To move each of these $2n$ distinct coefficients to the constant term's coefficient position in $V$ (i.e., shifting the coefficient $v_j$ to the leftmost term in $V$), the rotating computation of $V \cdot X^{-j}$ can use $2n$ distinct $j$ values, which are $\{0, 1, 2, \gap{$\cdots$}, n-1, n, \gap{$\cdots$}, 2n-1\}$, to move each of $(v_0, v_1, v_2, \gap{$\cdots$}, v_{n-1}, -v_0, -v_1, \gap{$\cdots$}, -v_{n-1})$ coefficients to the constant term's position. This implies that the exponent $j$ in $V\cdot X^{-j}$ can use any of the $2n$ distinct values to cover all possible $2n$ (sign-flipped) coefficient states of $V$. Also, remember that $j = \Delta m + e$.  Therefore, we will switch the modulo of $\Delta m + e$ from $q \rightarrow 2n$. Using the LWE modulus switching technique (\autoref{subsec:modulus-switch-lwe}), our original LWE ciphertext $\textsf{LWE}_{\vec{s}, \sigma}(\Delta m + e) = ({a}_0, {a}_1, \gap{$\cdots$} {a}_{k-1}, {b}) \in \mathbb{Z}_q^{k+1}$ (i.e., the initial input to the noise bootstrapping procedure) will be converted into the following: 

$ $

$\textsf{LWE}_{\vec{s}, \sigma}(\hat{\Delta} m) = (\hat{a}_0, \hat{a}_1, \gap{$\cdots$} \hat{a}_{k-1}, \hat{b}) \in \mathbb{Z}_{2n}^{k+1}$

$\vec{s} = (s_0, s_1, \gap{$\cdots$} s_{k-1}) \in \mathbb{Z}_2^{k}$  \textcolor{red}{ $\rhd$ the secret key stays the same, as each $s_i$ is binary}

$\hat{\Delta} = \Delta\dfrac{2n}{q} = \dfrac{2n}{t} \in \mathbb{Z}_{2n}$

$\hat{a}_i = \left \lceil a_i\dfrac{2n}{q} \right \rfloor \in \mathbb{Z}_{2n}$

$\hat e=\dfrac{2n}{q}e+\delta_b-\sum_{i=0}^{k-1}s_i\delta_{a_i}$, where $\delta_b=\hat b-\dfrac{2n}{q}b$ and $\delta_{a_i}=\hat a_i-\dfrac{2n}{q}a_i$ are component rounding errors. In particular,
\[
|\hat e|\leq\dfrac{2n}{q}|e|+\dfrac{1+\|\vec s\|_1}{2}.
\]
The simple integer-grid description assumes $t\mid 2n$ in addition to $t\mid q$.

$\hat{b} = \left \lceil b\dfrac{2n}{q} \right \rfloor \equiv \sum\limits_{i=0}^{k-1} \hat{a}_is_i + \hat{\Delta} m + \hat{e}  \in \mathbb{Z}_{2n}$



$ $

\para{The degree of $V(X)$ and Security:} The ring dimension is $n$, the polynomial degree is at most $n-1$, and the rotation period is $2n$. A lookup table needs room for noise-tolerant intervals, not just one coefficient per message. Ring dimension, modulus, secret distribution and error distribution must be selected jointly for security and correctness; increasing the degree alone is not a complete security argument.



\subsubsection{Halving the Plaintext Space To be Used}
\label{subsec:tfhe-zero-padding} 

In $\mathbb Z_q[X]/(X^n+1)$, the identity $X^n=-1$ forces a lookup table to be antiperiodic. Write $V(X)=\sum_{j=0}^{n-1}v_jX^j$, and let $r=j\bmod 2n$. Then
\[
[X^0]\bigl(V(X)X^{-j}\bigr)=F(j)=
\begin{cases}
v_r,&0\leq r<n,\\
-v_{r-n},&n\leq r<2n.
\end{cases}
\]
Here $[X^0]$ denotes the constant coefficient. Consequently $F(j+n)=-F(j)$; the second half of the table cannot be chosen independently of the first half. This is the actual restriction, rather than a prohibition on intermediate rotations or on negative rotation indices.

\para{Solution:} For an arbitrary lookup function, reserve one padding bit and restrict the admissible message set to a specified contiguous half of $\mathbb Z_t$ (with $t$ even). Design a constant lookup interval centered at each admissible $\hat\Delta m$, and leave a strict guard against rounding boundaries. A length-$n$ interval of phase indices can be assigned freely, and the other half is fixed by antiperiodicity. Signed messages may use negative phase indices; their coefficients must then include the sign dictated by $X^n=-1$.

\begin{tcolorbox}[title={\textbf{\tboxlabel{\ref*{subsec:tfhe-zero-padding}} Structure of Lookup Table Polynomial $\bm{V(X)}$}}]
Let $\hat\Delta=2n/t$, let $\mathcal D$ be the admissible half-domain, and choose an antiperiodic function $F:\mathbb Z_{2n}\to\mathbb Z_q$ such that
\[
F(\hat\Delta m+u)=m\quad(m\in\mathcal D,\ |u|<\hat\Delta/2),\qquad F(j+n)=-F(j).
\]
Boundary values use a fixed tie convention. Set
\[
V(X)=\sum_{j=0}^{n-1}F(j)X^j.
\]
The accumulator encrypts $\Delta V(X)$ at modulus $q$, so a successful lookup produces a constant term $\Delta m$, not an unscaled $m$. Correctness requires the total switched error (including component rounding) to stay in its assigned interval. For programmable bootstrapping replace $m$ on the right by the desired output $f(m)$ and choose the output scale explicitly.
\end{tcolorbox}

The encrypted lookup polynomial is also called an accumulator. A function already satisfying the required antiperiodicity may use a larger input domain; the restriction is on the function and its encoding.

\subsubsection{Blind Rotation}
\label{subsec:bootstrapping-blind-rotation} 

Blind rotation refers to rotating the coefficients of an \textit{encrypted} polynomial so that it is not possible to know how many positions the polynomial's coefficients have been rotated. After the rotation, it is not possible to see which coefficient has moved to which degree term. Blind rotation uses the basic polynomial rotation method 1 technique (Summary~\ref*{subsec:coeff-rotation} in \autoref{subsec:coeff-rotation}), with the difference that the $V(X)\cdot X^{-i}$ computation is done homomorphically. 

The phase convention in this chapter gives $X^{\hat\Delta m+\hat e_b}=X^{\hat b-\sum_{i=0}^{k-1}\hat a_i s_i}$. Therefore blind rotation starts with $X^{-\hat b}$ and applies conditional factors $X^{\hat a_i s_i}$. The evaluator is given $\textsf{BK}=\{\textsf{GGSW}_{\vec S_{bk},\sigma}^{\beta,l}(s_i)\}_{i=0}^{k-1}$, encrypting the input LWE secret bits under the output GLWE secret $\vec S_{bk}$. The evaluator does not receive that secret. Below $V_i$ denotes the unscaled lookup polynomial; its accumulator ciphertext encrypts $\Delta V_i$.

\begin{enumerate}
\item Initialize the GLWE accumulator encrypting $\Delta V$ under the output secret $\vec S_{bk}$, or use a trivial encryption for a public table.
\item Compute $V_0 = V \cdot X^{-\hat{b}}$, which is basically rotating $V$'s polynomials by $\hat{b}$ positions to the left. Since $\hat{b}$ is a known value visible in the LWE ciphertext, we can directly compute the rotation of $V_0 = V \cdot X^{-\hat{b}}$. 
\item Compute $V_1 = V_0 \cdot X^{\hat{a}_0s_0} = s_0 \cdot (V_0 \cdot X^{\hat{a}_0} - V_0) + V_0$.  
This formula works for the special case where $s_0 \in \{0, 1\}$: if $s_0 = 0$, then $V_1 = V_0$; else if $s_0 = 1$, then $V_1 = V_0 \cdot X^{\hat{a}_0}$. Computing $s_0 \cdot (V_0 \cdot X^{\hat{a}_0} - V_0) + V_0$ is done as a TFHE homomorphic addition and multiplication. We call this blind rotation of $V_0$, where the selection bit $s_0$ (i.e., the 1st element of the secret key vector $S$) is encrypted as a GGSW ciphertext by using $\vec{S}_{bk}$ (the output GLWE secret key) and $V_0$ is an encrypted polynomial as a GLWE ciphertext. Multiplying GLWE-encrypted $V_0$ with $x^{\hat{a}_0}$ is done by GLWE ciphertext-to-plaintext multiplication (\autoref{sec:glwe-mult-plain}), and subtracting the result by GLWE-encrypted $V_0$ is done by GLWE ciphertext-to-ciphertext addition/subtraction (\autoref{sec:glwe-add-cipher}), and multiplying the result by GGSW-encrypted $s_0$ is done by GLWE-to-GGSW multiplication (\autoref{subsubsec:tfhe-glwe-to-ggsw-multiplication}), and adding the result with GLWE-encrypted $V_0$ is done by GLWE ciphertext-to-ciphertext addition. If $s_0 = 1$, then the formula's $X^{\hat{a}_0}$ term gets multiplied to $V_0$, which effectively rotates $V_0$'s coefficients by $\hat{a}_0$ positions to the right. Else if $s_0 = 0$, then $V_0$ does not get rotated and stays the same. In both cases, the resulting $V_1$ is encrypted as a new GLWE ciphertext. The evaluator knows the public candidate rotation $\hat a_0$, but not the encrypted selection bit $s_0$. Hiding the selected branch relies on the security assumptions of the input and evaluation-key encryptions; possession of the appropriate secret key permits its recovery.
\item By using the same blind rotation method as in the previous step, compute the GLWE encryption of $V_2 = V_1 \cdot X^{\hat{a}_1s_1} = s_1 \cdot (V_1 \cdot X^{\hat{a}_1} - V_1) + V_1$. Note that we have the following publicly known components: $\hat{a}_1$, a GLWE encryption of $V_1$, and a GGSW ciphertext of $s_1$ encrypted by using $\vec{S}_{bk}$.
\item Continue to compute the GLWE encryption of $V_3, V_4, \gap{$\cdots$}, V_k$ in the same manner, and we finally get a GLWE encryption of $V_k$, whose computed value is equivalent to:

$V' = V_k$

$\text{ } = V_{k-1} \cdot X^{\hat{a}_{k-1}s_{k-1}}$

$\text{ } = V_0 \cdot X^{\hat{a}_0s_0}X^{\hat{a}_1s_1}\cdots X^{\hat{a}_{k-1}s_{k-1}}$

$\text{ } = V \cdot X^{- \hat{b} + \sum_{i=0}^{k-1}{\hat{a}_is_i}} $

$\text{ } = V \cdot X^{-(\hat{\Delta} m + \hat{e}_b)}$

\end{enumerate}

This means that the GLWE encryption of the final polynomial $V_k$ will have the coefficient $m$ in its constant term, as $V(X)$ is designed to have the mapping ($\hat\Delta m + \hat e_{*}, m$).

Intermediate rotations may cross any boundary. Only the final phase $\hat\Delta m+\hat e$ must belong to the admissible centered interval described in \autoref{subsec:tfhe-zero-padding}. This is a subset of $\mathbb Z_{2n}$, not $\mathbb Z_t$; it need not be $[0,n)$ when signed encodings are used.

\begin{figure}[h!]
    \centering
  \includegraphics[width=0.2\linewidth]{figures/mux.pdf}
  \caption{An illustration of the MUX logic gate }
  \label{fig:mux}
\end{figure}

\para{Homomorphic MUX Logic Gate:} In step 3, the formula $s_0 \cdot (V_0 \cdot X^{\hat{a}_0} - V_0) + V_0$ implements the MUX logic gate as shown in \autoref{fig:mux}, where in our case $s_0$ is the selection bit that chooses between the two inputs: $V_0$ and $V_0 \cdot X^{\hat{a}_0}$. If $s_0$ is 1, then the output is $V_0 \cdot X^{\hat{a}_0}$; otherwise, the output is $V_0$. In our design, the homomorphic computation of $s_0 \cdot (V_0 \cdot X^{\hat{a}_0} - V_0) + V_0$ effectively implements a homomorphic MUX gate, where the two data inputs are GLWE-encrypted and the selection bit is GGSW-encrypted.

\para{Dimensions of the GLWE and GGSW Ciphertexts:} A GLWE ciphertext contains $k'+1$ polynomials of degree at most $n-1$, hence $(k'+1)n$ coefficients. Each GGSW bootstrapping-key entry encrypts one input bit $s_i$ and contains $l(k'+1)^2$ polynomials, hence $l(k'+1)^2n$ coefficients. Flattening the output GLWE secret gives an LWE secret of dimension $k'n$. Choosing $k'=1$ yields RLWE/RGSW; this common efficiency choice still requires an appropriate security parameter set.

\subsubsection{Coefficient Extraction}

Next, use coefficient extraction (\autoref{subsec:tfhe-extraction}) on the rotated accumulator to obtain $\textsf{LWE}_{\vec{s}',\sigma}(\Delta m)$ under a vector $\vec{s}'$ of length $k'n$. Extraction itself adds no noise. If the input phase selects the correct lookup interval, the decoded message no longer depends on the old error $\hat e_b$; the output instead contains the accumulator and blind-rotation evaluation error $e_s$, which must satisfy its own decoding bound.

In fact, the homomorphic MUX logic in the blind rotation procedure (\autoref{subsec:bootstrapping-blind-rotation}) involves numerous ciphertext multiplications and additions, which can accumulate additional noise until we reach the point of coefficient extraction. To limit the accumulating noise during the series of MUX logic operations, we can carefully adjust the security parameters. For example, we can design a narrower Gaussian distribution for sampling the noise $e$, while designing a sufficiently large $n$ to compensate for the reduced noise. Meanwhile, increasing $q$ can make the system less secure because it becomes more vulnerable to lattice reduction attacks. 


\subsubsection{Key Switching}
\label{subsec:bootstrapping-key-switch} 

The output of the coefficient extraction is $\textsf{LWE}_{\vec{s}', \sigma}(\Delta m)$, which is encrypted by a secret key $\vec{s}'$ whose length is $k'\cdot n$. For consistency, we need to switch its key from $\vec{s}' \rightarrow \vec{s}$, where $\vec{s}$ is our original $k$-length key. This key switching can be done by using the technique explained in Summary~\ref*{subsec:tfhe-key-switching} (in \autoref{subsec:tfhe-key-switching}).



\subsubsection{Noise Bootstrapping Summary}
\label{subsec:tfhe-summary}

TFHE's noise bootstrapping procedure is summarized as follows:

\begin{tcolorbox}[title={\textbf{\tboxlabel{\ref*{subsec:tfhe-summary}} TFHE Noise Bootstrapping Procedure}}]

\textbf{\underline{Lookup Table Encryption}:} Initialize an encryption of $\Delta V(X)$ under the output GLWE secret $\vec S_{bk}$ of length $k'$, or use the trivial encryption for a public table. Blind rotation uses the public GGSW collection $\textsf{BK}$ defined above.

  
\begin{enumerate}
\item \textbf{\underline{Modulus Switch}:} Change the modulus of the TFHE ciphertext $\textsf{LWE}_{\vec{s}, \sigma}(\Delta m + e_b)$ from $q \rightarrow 2n$ to get $\textsf{LWE}_{\vec{s}, \sigma}(\hat{\Delta} m + \hat e_b)$, where $\hat\Delta = \Delta \cdot \dfrac{2n}{q} = \dfrac{2n}{t}$.

$ $

\item \textbf{\underline{Blind Rotation}:} Rotate the GLWE-encrypted polynomial $V$ by $(\hat{b} - \sum\limits_{i=0}^{k-1}{\hat{a}_is_i}) = (\hat{\Delta}m + \hat{e}_b)$ positions to the left, by recursively computing: 

$V_0 = V \cdot X^{-\hat{b}}$

$V_1 = V_0 \cdot X^{\hat{a}_0s_0} = s_0 \cdot (V_0 \cdot X^{\hat{a}_0} - V_0) + V_0$

\text{ } $\vdots$

$V_k = V_{k-1} \cdot X^{\hat{a}_{k-1}s_{k-1}} = s_{k-1} \cdot (V_{k-1} \cdot X^{\hat{a}_{k-1}} - V_{k-1}) + V_{k-1}$

$= V_0 \cdot X^{\hat{a}_0s_0}X^{\hat{a}_1s_1}\cdots X^{\hat{a}_{k-1}s_{k-1}}$

$= V \cdot X^{- \hat{b} + \sum_{i=0}^{k-1}{\hat{a}_is_i}} $

$ = V \cdot X^{-(\hat{\Delta}m + \hat{e}_b)}$

$ $

\noindent Each step of the actual blind rotation above is computed as the following TFHE ciphertext-to-ciphertext multiplication and addition:

$\textsf{GLWE}_{\vec{S}_{bk}, \sigma}(\Delta V_{i+1}) = \textsf{GGSW}_{\vec{S}_{bk}, \sigma}^{\beta, l}(s_i) \cdot (\textsf{GLWE}_{\vec{S}_{bk}, \sigma}(\Delta V_i) \cdot X^{\hat{a}_i} - \textsf{GLWE}_{\vec{S}_{bk}, \sigma}(\Delta V_i)) + \textsf{GLWE}_{\vec{S}_{bk}, \sigma}(\Delta V_i)$

$ $

\item \textbf{\underline{Coefficient Extraction}:} Homomorphically extract the constant term's coefficient $m$ from the rotated polynomial $V_k$, which is $\textsf{LWE}_{\vec{s}', \sigma}(\Delta m)$, where $\vec{s}'$ is a vector of length $k' \cdot n$. 

$ $

\item \textbf{\underline{Key Switching}:} Homomorphically switch the key of the LWE ciphertext from $\vec{s}' \rightarrow \vec{s}$. 

%\item \textbf{Rescaling:} Multiply $\textsf{LWE}_{\vec{s}, \sigma}(\hat{\Delta}m)$ by $\dfrac{q}{2n}$ to get:  

%$\dfrac{q}{2n} \cdot LWE_{S, \sigma}(\hat{\Delta}m)$

%$\textsf{LWE}_{\vec{s}, \sigma}(\dfrac{q}{2n} \cdot \hat{\Delta}m)$

%$ = LWE_{S, \sigma}(\Delta m)$

%, whose noise get reinitialized to $e_s$. 
\end{enumerate}

$ $

\para{\underline{Halving the Usable Plaintext Range}:} Reserve a specified half-domain for an arbitrary LUT and center each lookup interval at its encoded message. The table must satisfy $F(j+n)=-F(j)$, and the full switched error must stay inside the assigned interval. Signed phase indices and intermediate rotations beyond $n-1$ are permitted.

$ $

\para{\underline{Choice of $k'$}:} For computational efficiency, $k'$ is set to be 1, which simplifies the GLWE and GGSW ciphertexts as RLWE and RGSW ciphertexts.  
%\item Piping the LUT output to a \textbf{negacyclic function}.
%\item Multiplying the LUT output with \textbf{another negacyclically rotating polynomial}. 
%\end{itemize}
\end{tcolorbox}


\subsubsection{Example: Noise Bootstrapping}
\label{subsec:tfhe-noise-bootstrapping-ex}

Suppose the GLWE security setup: $n = 16$, $t = 8$, $q = 64$, $k = 8$

$\mathbb{Z}_{t=8} = \{-4, -3, -2, -1, 0 , 1, 2, 3\}$

$\mathbb{Z}_{q=64} = \{ -32, -31, -30, \gap{$\cdots$}, 29, 30, 31 \}$

$\Delta = \dfrac{q}{t} = \dfrac{64}{8} = 8$

$ $

And suppose we have the following LWE ciphertext:

$\vec{s} = (1, 0, 0, 1, 1, 1, 0, 1) \in \mathbb{Z}_2^{k=8}$

$m = 1  \in \mathbb{Z}_{t=8}$

$\Delta m = 1 \cdot 8 = 8 \in \mathbb{Z}_{q=64}$



$\textsf{LWE}_{\vec{s}, \sigma}(\Delta m) = (a_0, a_1, a_2, a_3, a_4, a_5, a_6, a_7, b) 
= (8, -28, 4, -32, 0, 31, -6, 7, 24) \in \mathbb{Z}_{q=64}^{k+1=9}$

$e = 2 \in \mathbb{Z}_{q=64}$ (should be the case that $|e| < \dfrac{\Delta}{2} = 4$ for correct decryption)

$b = \sum\limits_{i=0}^{7}a_is_i + \Delta m + e = (8 -32 + 31 +7) + 8 + 2 = 24 \in \mathbb{Z}_{q=64}$

$ $

Now then, the TFHE noise bootstrapping procedure is as follows:

$ $

\begin{enumerate}
\item \textbf{Modulus Switch:} Switch the modulus of $\textsf{LWE}_{\vec{s}, \sigma}(\Delta m)$ From $q \rightarrow 2n$, which is from $64 \rightarrow 32$. After the modulus switch, the original LWE ciphertext is converted as follows:

$\mathbb{Z}_{2n=32} = \{-16, -15, -14, \gap{$\cdots$}, 13, 14, 15\}$

$\vec{s} = (1, 0, 0, 1, 1, 1, 0, 1) \in \mathbb{Z}_2^{k=8}$


$\hat{\Delta} = \Delta\dfrac{2n}{64} = 8\dfrac{32}{64} = 4$

$\hat{\Delta}m = 4 \cdot 1 = 4 \in \mathbb{Z}_{2n=32}$

$\hat e=\hat b-\sum_{i=0}^{7}\hat a_i s_i-\hat\Delta m=12-(4-16+16+4)-4=0\pmod{32}$. The component $\hat a_5=16$ is a noncentered representative of $-16$; either represents the same residue.

$ $

$\textsf{LWE}_{\vec{s}, \sigma}(\hat{\Delta}m) = (\hat{a}_0, \hat{a}_1, \hat{a}_2, \hat{a}_3, \hat{a}_4, \hat{a}_5, \hat{a}_6, \hat{a}_7, \hat{b}) \in \mathbb{Z}_{2n=32}^{k+1=9}$


$ = (\Big\lceil 8\dfrac{32}{64} \Big\rfloor, \Big\lceil -28\dfrac{32}{64}\Big\rfloor, \Big\lceil 4\dfrac{32}{64}\Big\rfloor, \Big\lceil -32\dfrac{32}{64}\Big\rfloor, \Big\lceil 0\dfrac{32}{64}\Big\rfloor, \Big\lceil 31\dfrac{32}{64}\Big\rfloor, \Big\lceil -6\dfrac{32}{64}\Big\rfloor, \Big\lceil 7\dfrac{32}{64}\Big\rfloor, \Big\lceil 24\dfrac{32}{64}\Big\rfloor)$

$ $

$= (4, {-14}, 2, -16, 0, 16, {-3}, 4, 12)$

$ $

The exact switched phase is $12-(4-16+16+4)=4\pmod{32}$. The scaled input error $2\cdot32/64=1$ is canceled here by the component-rounding contribution $-1$, giving $\hat e=0$. In general it is incorrect to replace the total switched error by the rounded value of $2ne/q$.

\item \textbf{Blind Rotation:} Use the admissible messages $\{-2,-1,0,1\}$ and $|\hat e|<2$. The table is centered at the phase values $-8,-4,0,4$. At the unused tie boundaries choose rounding toward the higher integer. The unscaled lookup coefficients are
\[
(v_0,\ldots,v_{15})=(0,0,1,1,1,1,2,2,2,2,1,1,1,1,0,0).
\]
\begin{table}[h]
\centering
\begin{tabular}{c|c|c|c|c}
\hline
Phase $j$ in $V X^{-j}$ & $-10,-9,-8,-7$ & $-6,-5,-4,-3$ & $-2,-1,0,1$ & $2,3,4,5$\\
\hline
Constant coefficient & $-2$ & $-1$ & $0$ & $1$\\
\hline
\end{tabular}
\caption{Centered lookup for $n=16,q=64,t=8$. Boundary points are assigned a tie convention; correctness uses the strict bound $|\hat e|<2$.}
\label{tab:lut}
\end{table}
Negative indices select the negated coefficients dictated by $X^{16}=-1$. The accumulator encrypts $\Delta V=8V$; the lookup indices use the different scale $\hat\Delta=4$.

During the actual bootstrapping, we will do a blind rotation of \autoref{tab:lut}'s $V(X)$ (which is GLWE-encrypted) by $\hat{b} - \sum\limits_{i=0}^{7}\hat{a}_is_i = 4$ positions to the left, which is computed as follows:

$\hat\Delta m + \hat e = \hat{b} - \sum\limits_{i=0}^{7}\hat{a}_is_i = 12 - (4 - 16 + 16 + 4) = 4 \text{ mod 32} \in \mathbb{Z}_{2n=32}$ 


In \autoref{tab:lut}, phase $4$ selects $v_4=1=m$. The accumulator contains $\Delta v_4=8$, so extraction produces a ciphertext decrypting to $m=1$, not the scalar ciphertext ``$1$''. The toy example uses $k'=8$ output GLWE polynomials, hence extraction gives a $128$-dimensional LWE secret before key switching back to dimension $8$; these small parameters are illustrative, not secure.



$ $

The actual blind rotation is computed as follows:

$\vec{s} = (1, 0, 0, 1, 1, 1, 0, 1)$

$\textsf{LWE}_{\vec{s}, \sigma}(\hat{\Delta}m) = (\hat{a}_0, \hat{a}_1, \hat{a}_2, \hat{a}_3, \hat{a}_4, \hat{a}_5, \hat{a}_6, \hat{a}_7, \hat{b}) =  (4, -14, 2, -16, 0, 16, -3, 4, 12)$

$V_0 = V \cdot X^{-\hat{b}} = V \cdot X^{-12} = v_{12} + v_{13}X + v_{14}X^2 + \cdots $

$V_1 = V_0 \cdot X^{\hat{a}_0s_0} = s_0 \cdot (V_0 \cdot X^{\hat{a}_0} - V_0) + V_0 = V_0 \cdot X^{4} = v_{8} + v_{9}X + v_{10}X^2 + \cdots$ 

$V_2 = V_1 \cdot X^{\hat{a}_1s_1} = s_1 \cdot (V_1 \cdot X^{\hat{a}_1} - V_1) + V_1 = V_1 = v_{8} + v_{9}X + v_{10}X^2 + \cdots$

$V_3 = V_2 \cdot X^{\hat{a}_2s_2} = s_2 \cdot (V_2 \cdot X^{\hat{a}_2} - V_2) + V_2 = V_2 = v_{8} + v_{9}X + v_{10}X^2 + \cdots$

$V_4 = V_3 \cdot X^{\hat{a}_3s_3} = s_3 \cdot (V_3 \cdot X^{\hat{a}_3} - V_3) + V_3 = V_3 \cdot X^{-16} = -v_{8} - v_{9}X - v_{10}X^2 - \cdots$

$V_5 = V_4 \cdot X^{\hat{a}_4s_4} = s_4 \cdot (V_4 \cdot X^{\hat{a}_4} - V_4) + V_4 = V_4 \cdot X^{0} = -v_{8} - v_{9}X - v_{10}X^2 - \cdots$

$V_6 = V_5 \cdot X^{\hat{a}_5s_5} = s_5 \cdot (V_5 \cdot X^{\hat{a}_5} - V_5) + V_5 = V_5 \cdot X^{16} = v_{8} + v_{9}X + v_{10}X^2 + \cdots$

$V_7 = V_6 \cdot X^{\hat{a}_6s_6} = s_6 \cdot (V_6 \cdot X^{\hat{a}_6} - V_6) + V_6 = V_6 = v_{8} + v_{9}X + v_{10}X^2 + \cdots$

$V_8 = V_7 \cdot X^{\hat{a}_7s_7} = s_7 \cdot (V_7 \cdot X^{\hat{a}_7} - V_7) + V_7 = V_7 \cdot X^{4} = v_{4} + v_5X + v_6X^2 + \cdots $

$ $

The final output of blind rotation is the GLWE ciphertext of $V_8$, $\textsf{GLWE}_{\vec{S}, \sigma}(\Delta V_8)$, whose constant term's coefficient is $v_4 = m = 1$.

$ $

Each step of the actual blind rotation above is computed as the following TFHE ciphertext-to-ciphertext multiplication:

$\textsf{GLWE}_{\vec{S}, \sigma}(\Delta V_{i+1}) = \textsf{GGSW}_{\vec{S}, \sigma}^{\beta, l}(s_i) \cdot (\textsf{GLWE}_{\vec{S}, \sigma}(\Delta V_i) \cdot X^{\hat{a}_i} - \textsf{GLWE}_{\vec{S}, \sigma}(\Delta V_i)) + \textsf{GLWE}_{\vec{S}, \sigma}(\Delta V_i)$

$ $

We will leave this computation for the reader's exercise.

$ $

\item \textbf{Coefficient Extraction:} At the end of blind rotation, we finally get the following GLWE ciphertext:

$\textsf{GLWE}_{\vec{S}, \sigma}(\Delta V_8)$

$ = \textsf{GLWE}_{\vec{S}, \sigma}\bm(\Delta \cdot (v_4 + v_5X + v_6X^2 + v_7X^3 + v_8X^{4} + v_9X^{5} + v_{10}X^{6} + v_{11}X^{7} + v_{12}X^{8} + v_{13}X^{9} + v_{14}X^{10} + v_{15}X^{11} - v_{0}X^{12} - v_{1}X^{13} - v_{2}X^{14} - v_{3}X^{15})\bm)$

$ = \textsf{GLWE}_{\vec{S}, \sigma}\bm(\Delta\cdot(1+X+2X^2+2X^3+2X^4+2X^5+X^6+X^7+X^8+X^9-X^{14}-X^{15})\bm)$

$ = \left(A_0 = \sum\limits_{j=0}^{15}a_{0,j}X^j, A_1 = \cdots, A_{k-1} = \cdots, B = \sum\limits_{j=0}^{15}b_{j}X^j\right)$

Now, we extract the constant term's coefficient of the encrypted polynomial $\textsf{GLWE}_{\vec{S}, \sigma}(\Delta \cdot (1+X+2X^2+\cdots))$ by using the coefficient extraction formula (Summary~\ref{subsec:tfhe-extraction}). Specifically, we will extract the constant term's coefficient, which corresponds to $\textsf{LWE}_{\vec{s}\,\prime, \sigma}(\Delta m_0)$. We extract $\textsf{LWE}_{\vec{s}\,\prime, \sigma}(\Delta m_0)$ by computing the following:

$\textsf{LWE}_{\vec{s}\,\prime, \sigma}(\Delta m_0) = (a_0', a_1', \gap{$\cdots$} , a_{nk-1}', b_h)$ \textcolor{red}{ $\rhd$ where $h = 0$}

\[
    \text{, where } a'_{n \cdot i + j} =   
    \begin{cases}
      a_{i,0 - j} \text{ (if } 0 \leq j \leq 0\text{)}\\
      -a_{i,n + 0 - j} \text{ (if } 0+1 \leq j \leq n-1\text{)}\\
  \end{cases}
, b_0 \text{ is obtained from the polynomial } B
\]


Finally use \autoref{subsec:tfhe-key-switching} to switch the extracted ciphertext from the flattened key $\vec s\,\prime$ back to the original $\vec s$.

\end{enumerate}


\subsubsection{Discussion}
\label{subsec:tfhe-noise-bootstrapping-discussion}

\begin{itemize}
\item \textbf{Programmable Bootstrapping}: Replace the identity lookup in \autoref{subsec:tfhe-noise-bootstrapping} by a table for $f(m)$, with centered input intervals, an explicit output scale and the required negacyclic layout. This yields gates such as those in \autoref{subsec:tfhe-noise-bootstrapping-gate}. An encrypted table can hide table entries under its encryption assumptions, but this alone is not a construction of functional encryption with function-restricted decryption keys. Identity and programmable bootstrapping have comparable noise bounds for the same accumulator and evaluation parameters; an encrypted table's initial noise and the chosen encoding must also be counted.
\item \textbf{Bootstrapping Noise}: The CMux operations in \autoref{subsec:bootstrapping-blind-rotation} introduce evaluation-key and decomposition errors. These form part of the new output noise; coefficient extraction does not erase them or introduce additional noise. Correctness needs both a successful table selection and a bound on the accumulated output error.
\end{itemize}


\subsubsection{Application: Gate Bootstrapping}
\label{subsec:tfhe-noise-bootstrapping-gate}

In this example $n=8$. The error bound below concerns the total phase error after forming the linear combination and modulus switching. The accumulator is an encryption of $\Delta V$ at modulus $q$, so the extracted output uses the same message scale as the input.

Besides the CMux used during blind rotation (\autoref{subsec:bootstrapping-blind-rotation}), a LUT can implement Boolean gates such as AND, NAND, OR and XOR. In this example $q=32$, $t=8$, $n=8$, $\Delta=q/t=4$ and $\hat\Delta=2n/t=2$. Encode Boolean inputs by $0\mapsto-1$ and $1\mapsto1$. Assume the total error after the affine combination and modulus switching satisfies $\hat e\in[-1,1]$. The three possible affine phases are then centered at $-6,-2,2$.


\begin{table}[ht]
  \centering
  \resizebox{\columnwidth}{!}{%
    \begin{tabular}{c c | r r | c | c | c | c}
      \hline
        \multicolumn{8}{c}{Lookup Table Polynomial $V(X) = 1 + 1X + 1X^2 + 1X^3 + 1X^4 + 1X^5 + 1X^6 + 1X^7$}\\
      \hline
      \multicolumn{2}{c|}{\textbf{Decoded}} & 
      \multicolumn{2}{c|}{\textbf{Encoded}} & 
      \multicolumn{1}{c|}{\begin{tabular}{c}
        \textbf{Linear} \\ \textbf{combination} \\
        \textbf{of encodings}
      \end{tabular}} & 
      \multicolumn{1}{c|}{\begin{tabular}{c}
        \textbf{Scaled} \textbf{Encoded} \\
         \textbf{Combination} 
      \end{tabular}} &
      \textbf{Bootstrapping} & 
      \begin{tabular}{c}
        \textbf{Decoded} \\ \textbf{result} \\
      \end{tabular} \\
        \hline
      $d_1$ & $d_2$ & $m_1$ & $m_2$ & $m_1 + m_2 - 1$ & $\hat\Delta\cdot(m_1+m_2-1)$ &
      $V(X)\cdot X^{-(\hat\Delta\cdot(m_1 + m_2 - 1)+\hat e)}$ &   $d_1 \land d_2$ \\
      \hline
      0 & 0 &
      $-1$ & $-1$ &
      $-3$ &
      $-6$ &
      constant term's coeff. is $-1$ &
      0 \\
      0 & 1 & 
      $-1$ & $1$ &
      $-1$ &
      $-2$ &
      constant term's coeff. is $-1$ &
      0 \\
      1 & 0 & 
      $1$ & $-1$ &
      $-1$ &
      $-2$ &
      constant term's coeff. is $-1$ &
      0 \\
      1 & 1 & 
      $1$ & $1$ &
      $1$ &
      $2$ &
      constant term's coeff. is $1$ &
      1 \\
      \hline
    \end{tabular}
  }
  \caption{An example truth table for an AND operation with an additional encoding.}
  \label{tab:gate-and}
\end{table}


%\begin{figure}[h!]
%    \centering
%  \includegraphics[width=\linewidth]{figures/torus-gate.png}
%  \caption{An illustration of an AND gate in Torus \href{https://assets-global.website-files.com/622ef9de9152c97467eac748/6298a348a0c75850f0688c83_NAND_bootstrap.png}{(Source)}}
%  \label{fig:torus-gate}
%\end{figure}


In \autoref{tab:gate-and}, $V(X)=\sum_{i=0}^7X^i$ returns $-1$ at phases $-6,-2$ and $+1$ at phase $2$, including the allowed error intervals. Form the affine ciphertext as $\textsf{ct}_1+\textsf{ct}_2-(0,\Delta)$; subtracting an unscaled $1$ would encode a different offset. Initialize the accumulator with $\Delta V$. Extraction and the usual key switch then produce an LWE encryption at scale $\Delta$ of the encoded AND result. The error bound is on the aggregate switched phase, not on each input error separately; larger supported ranges require a new correctness and security parameter analysis.

This gate table is antiperiodic, so it may use both halves of the $2n=16$ rotation indices consistently. For a general lookup function, reserve the admissible half-domain and centered intervals specified in \autoref{subsec:tfhe-zero-padding}. The identity lookup of \autoref{subsec:bootstrapping-overview} likewise needs a compatible encoding; the restriction is not a ban on rotations exceeding $n-1$.

Besides the AND gate, other logic gates can be built in a similar manner, each of which is based on a different linear combination formula and LUT layout. 

\para{Division:} Modular inversion, integer division and rounding are different operations. A nonunit has no multiplicative inverse in $\mathbb Z_q$, but this does not prevent a lookup table from evaluating division by a public constant or another bounded unary division function on its specified domain. General integer division can also be implemented with Boolean circuits and gate bootstrapping; division by zero needs an application-defined convention.

\subsubsection{Application: Neural Networks Bootstrapping}
\label{subsec:tfhe-neural-network}

\begin{figure}[h!]
    \centering
  \includegraphics[width=0.8\linewidth]{figures/neural-network.pdf}
  \caption{An illustration of neural networks }
  \label{fig:neural-network}
\end{figure}

\begin{figure}[h!]
    \centering
  \includegraphics[width=0.8\linewidth]{figures/nn-homomorphic.pdf}
  \caption{Programmable bootstrapping in a quantized neural network. The diagram's added bias denotes the encoded constant $(0,\Delta b)$; the text specifies the quantization, scales and admissible lookup ranges.}
  \label{fig:neural-network2}
\end{figure}

Homomorphic encryption can be applied to the neurons of deep neural networks, in which each neuron is generally comprised of two steps of computation: 

\begin{enumerate}
  \item \textbf{Linear Combination of Input Values:} An input feature value (or intermediate value) set $(x_1, x_2, \gap{$\cdots$}, x_n)$, a weight set $(w_1, w_2, \gap{$\cdots$}, w_n)$, and a bias $b$ are computed as: $y = \sum\limits_{i=1}^{n}x_iw_i + b$.
\item \textbf{Activation Function:} $f(y)$ is computed, where $f$ is a non-linear activation function such as the $\sin$ function, ReLU, sigmoid, hyperbolic tangent, etc. 
\end{enumerate}

With quantized integer features and public integer weights, the first step is implemented as
\[
\sum_{i=1}^n\textsf{LWE}_{\vec s,\sigma}(\Delta x_i)w_i+(0,\Delta b)
\]
. It uses ciphertext addition (\autoref{sec:glwe-add-cipher}) and ciphertext-to-plaintext multiplication (\autoref{sec:glwe-mult-plain}). The input and bias scales must agree, and the weighted sum's range and noise must fit the next lookup.

The activation can be evaluated through a finite lookup table after choosing input and output quantization, clipping ranges and scales. Nonpolynomial real functions cannot be represented exactly by a finite real arithmetic polynomial, although polynomial approximation and interpolation on finite domains are possible. Programmable bootstrapping evaluates the selected discrete function and refreshes its ciphertext. Repeating this supports deep networks only when each linear stage stays within the lookup domain and noise budget, and the accumulated failure probability is acceptable.



\subsection{TFHE on a Discrete Torus}
\label{subsec:torus}

\textbf{- Reference:} 
\href{https://eprint.iacr.org/2021/1402.pdf}{Guide to
Fully Homomorphic Encryption
over the [Discretized] Torus}~\cite{torus}

$ $

Torus $\mathbb{T}$ is a continuous real number domain between 0 and 1 that wraps around, that is $[0, 1)$. 

A discrete torus $\mathbb{T}_t$ is a finite real number set: $\Big(0, \dfrac{1}{t}, \dfrac{2}{t}, \gap{$\cdots$}, \dfrac{t - 1}{t}\Big)$


In the previous subsections, we learned the TFHE scheme based on the following setup:

$ $

$m \in \mathbb{Z}_t$ 
$\vec{s} \in \{0, 1\}^{k}$
$e \in \mathbb{Z}_q$

$\textsf{LWE}_{\vec{s}, \sigma}(\Delta m) = (a_0, a_1, \gap{$\cdots$}, a_{k-1}, b) \in \mathbb{Z}_q^{k + 1}$

$ $

Dividing the integer ciphertext components by $q$ gives the corresponding discrete-torus representation (the original theory can also use the continuous torus):


$ $

$m \in \mathbb{T}_{t}, \text{ } \vec{s} \in \{0, 1\}^{k}, \text{ } e \in \mathbb{T}_{q}$

$ $

$\textsf{LWE}_{\vec{s}, \sigma}(m) = \textsf{ct} = (a_0, a_1, \gap{$\cdots$}, a_{k-1}, b) \in \mathbb{T}_{q}^{k+1}$

$b = \sum\limits_{i=0}^{k-1}(a_is_i) + m + e \in \mathbb{T}_q$

$\textsf{LWE}_{\vec{s}, \sigma}^{-1}(\textsf{ct}) =  \Big\lceil b - \sum\limits_{i=0}^{k-1}(a_is_i) \Big\rfloor_{\frac{1}{t}} = \Big\lceil m + e \Big\rfloor_{\frac{1}{t}} = m$, given $|e| < \dfrac{1}{2t}$

\textcolor{red}{ $\rhd$ where $\lceil x \rfloor_{\frac{1}{t}}$ means rounding $x$ to the nearest multiple of $\dfrac{1}{t}$}

$ $

The torus is the additive group $\mathbb T=\mathbb R/\mathbb Z$, with multiplication by integers; arbitrary real multiplication is not a well-defined ring operation on it. A discrete implementation represents $a/q\bmod1$ by the integer $a\bmod q$ (for example with fixed-width integers), so floating-point storage is not required. Normalizing the integer encoding by $q$ sends $\Delta m_{\rm int}$ to $m_{\rm int}/t$ when $t\mid q$. Thus the embedding scale and message spacing $1/t$ remain present even when the symbol $\Delta$ is omitted. Correctness requires the centered torus error to satisfy $|e|<1/(2t)$. FFT arithmetic may use floating point internally; that is separate from the ciphertext representation.

Both the torus-based and integer-ring-based TFHE schemes are built based on the same fundamental principles.